% ===========================================================================
% TESIS DE GRADO — VulnSeeker
% Universidad del Zulia — Facultad Experimental de Ciencias
% Licenciatura en Computación
% Autor: Claudio Enrique Palmar León
% Tutor: Sigerist Rodríguez
% Formato: APA 7ma edición
% ===========================================================================
\documentclass[12pt, a4paper]{article}

% ── Codificación y lenguaje ────────────────────────────────────────────────
\usepackage[utf8]{inputenc}
\usepackage[spanish, es-nodecimaldot]{babel}

% ── Fuente Helvetica ──────────────────────────────────────────────────────
\usepackage[scaled]{helvet}
\renewcommand{\familydefault}{\sfdefault}
\usepackage[T1]{fontenc}

% ── Márgenes APA (2.54 cm = 1 pulgada) ──────────────────────────────────
\usepackage{geometry}
\geometry{
    top=2.54cm,
    bottom=2.54cm,
    left=2.54cm,
    right=2.54cm,
    headheight=15pt
}

% ── Interlineado doble (APA 7) ───────────────────────────────────────────
\usepackage{setspace}
\doublespacing

% ── Sangría de primera línea (APA 7: 1.27 cm) ───────────────────────────
\usepackage{indentfirst}
\setlength{\parindent}{1.27cm}

% ── Paquetes generales ───────────────────────────────────────────────────
\usepackage{graphicx}
\usepackage{booktabs}
\usepackage{float}
\usepackage{longtable}
\usepackage{array}
\usepackage{enumitem}
\usepackage{amsmath}
\usepackage{amssymb}
\usepackage{url}
\usepackage{xcolor}
\usepackage{listings}
\usepackage{csquotes}
\usepackage{caption}

% ── Bibliografía APA 7 ──────────────────────────────────────────────────
\usepackage[style=apa, backend=biber, language=spanish]{biblatex}
\DeclareLanguageMapping{spanish}{spanish-apa}
\addbibresource{referencias.bib}

% ── Hipervínculos ────────────────────────────────────────────────────────
\usepackage[colorlinks=true, linkcolor=black, citecolor=black, urlcolor=blue]{hyperref}

% ── Encabezados (running head APA) ──────────────────────────────────────
\usepackage{fancyhdr}
\pagestyle{fancy}
\fancyhf{}
\fancyhead[L]{\small\textit{VULNSEEKER: ESCÁNER MODULAR DE VULNERABILIDADES}}
\fancyhead[R]{\thepage}
\renewcommand{\headrulewidth}{0pt}

% ── Configuración de código fuente ──────────────────────────────────────
\lstset{
    basicstyle=\footnotesize\ttfamily,
    breaklines=true,
    frame=single,
    numbers=left,
    numberstyle=\tiny,
    tabsize=4
}

% ===========================================================================
\begin{document}

% ===========================================================================
% PORTADA
% ===========================================================================
\begin{titlepage}
    \centering
    \singlespacing

    \vspace*{1cm}

    {\footnotesize
        REPÚBLICA BOLIVARIANA DE VENEZUELA\\
        UNIVERSIDAD DEL ZULIA\\
        FACULTAD EXPERIMENTAL DE CIENCIAS\\
        DIVISIÓN DE ESTUDIOS BÁSICOS SECTORIALES\\
        LICENCIATURA EN COMPUTACIÓN
    }

    \vspace{4cm}

    {\LARGE \textbf{ESCÁNER MODULAR PARA LA DETECCIÓN\\
    AUTOMATIZADA DE VULNERABILIDADES\\
    EN APLICACIONES WEB}}

    \vspace{1cm}

    {\large Trabajo presentado como requisito para optar por el título de\\
    Licenciado en Computación}

    \vspace{4cm}

    \begin{flushright}
        \textbf{Autor:} Claudio Enrique Palmar León\\
        \textbf{C.I.:} V-30.106.822\\[0.5cm]
        \textbf{Tutor:} Prof. Sigerist Rodríguez
    \end{flushright}

    \vfill

    {\large Maracaibo, 2026}
\end{titlepage}

% ===========================================================================
% PÁGINAS PRELIMINARES
% ===========================================================================
\pagenumbering{roman}

% ── Dedicatoria (placeholder) ────────────────────────────────────────────
\newpage
\thispagestyle{empty}
\vspace*{5cm}
\begin{flushright}
    \textit{[Dedicatoria]}
\end{flushright}

% ── Agradecimientos (placeholder) ────────────────────────────────────────
\newpage
\thispagestyle{empty}
\section*{Agradecimientos}
\addcontentsline{toc}{section}{Agradecimientos}
\textit{[Por redactar]}

% ── Resumen ──────────────────────────────────────────────────────────────
\newpage
\section*{Resumen}
\addcontentsline{toc}{section}{Resumen}

\textbf{Palmar León, Claudio Enrique.} \textit{Escáner Modular para la Detección
Automatizada de Vulnerabilidades en Aplicaciones Web.} Trabajo de Grado.
Universidad del Zulia. Facultad Experimental de Ciencias. Licenciatura en
Computación. Maracaibo, Venezuela. 2026.

\vspace{0.5cm}

La presente investigación tuvo como objetivo desarrollar un escáner modular para
la detección automatizada de vulnerabilidades en aplicaciones web, denominado
VulnSeeker. A partir de un diagnóstico del estado del arte en herramientas de
evaluación dinámica de seguridad (DAST), se diseñó una arquitectura extensible en
Python compuesta por un motor de orquestación, un crawler autenticado, un sistema
de fingerprinting tecnológico y 25~módulos de detección que cubren 8 de las
10~categorías del OWASP Top~10. La herramienta incorpora consultas en tiempo real
a la base de datos nacional de vulnerabilidades del NIST (NVD~API~2.0) y genera
reportes técnicos con análisis asistido por inteligencia artificial. La validación
se realizó mediante 150~pruebas unitarias automatizadas y escaneos controlados
sobre entornos de práctica estándar. Los resultados demostraron que la integración
de múltiples motores de detección en una arquitectura modular incrementa la
cobertura de vulnerabilidades detectadas sin sacrificar extensibilidad. VulnSeeker
se distribuye como software de código abierto, contribuyendo a la democratización
del acceso a herramientas de auditoría de seguridad web.

\vspace{0.5cm}
\textbf{Palabras clave:} seguridad web, DAST, OWASP, vulnerabilidades, escáner
modular, código abierto, Python.

% ── Índice general ──────────────────────────────────────────────────────
\newpage
\tableofcontents

% ── Lista de tablas ─────────────────────────────────────────────────────
\newpage
\listoftables

% ── Lista de figuras ────────────────────────────────────────────────────
\newpage
\listoffigures

% ── Lista de Abreviaturas ──────────────────────────────────────────────
\newpage
\section*{Lista de Abreviaturas}
\addcontentsline{toc}{section}{Lista de Abreviaturas}

\begin{tabular}{p{3cm} p{10cm}}
    \textbf{API}   & Application Programming Interface \\
    \textbf{CI/CD} & Continuous Integration / Continuous Deployment \\
    \textbf{CLI}   & Command-Line Interface \\
    \textbf{CMS}   & Content Management System \\
    \textbf{CORS}  & Cross-Origin Resource Sharing \\
    \textbf{CSRF}  & Cross-Site Request Forgery \\
    \textbf{CVE}   & Common Vulnerabilities and Exposures \\
    \textbf{CVSS}  & Common Vulnerability Scoring System \\
    \textbf{CWE}   & Common Weakness Enumeration \\
    \textbf{DAST}  & Dynamic Application Security Testing \\
    \textbf{GUI}   & Graphical User Interface \\
    \textbf{HSTS}  & HTTP Strict Transport Security \\
    \textbf{HTTP}  & HyperText Transfer Protocol \\
    \textbf{HTTPS} & HyperText Transfer Protocol Secure \\
    \textbf{IA}    & Inteligencia Artificial \\
    \textbf{IAST}  & Interactive Application Security Testing \\
    \textbf{LFI}   & Local File Inclusion \\
    \textbf{LLM}   & Large Language Model \\
    \textbf{NVD}   & National Vulnerability Database \\
    \textbf{NIST}  & National Institute of Standards and Technology \\
    \textbf{OSINT} & Open Source Intelligence \\
    \textbf{OWASP} & Open Worldwide Application Security Project \\
    \textbf{PYME}  & Pequeña y Mediana Empresa \\
    \textbf{RFI}   & Remote File Inclusion \\
    \textbf{SAST}  & Static Application Security Testing \\
    \textbf{SQL}   & Structured Query Language \\
    \textbf{SQLi}  & SQL Injection \\
    \textbf{SSRF}  & Server-Side Request Forgery \\
    \textbf{TLS}   & Transport Layer Security \\
    \textbf{URL}   & Uniform Resource Locator \\
    \textbf{UX}    & User Experience \\
    \textbf{XSS}   & Cross-Site Scripting \\
\end{tabular}

% ── Glosario de Términos ──────────────────────────────────────────────
\newpage
\section*{Glosario de Términos}
\addcontentsline{toc}{section}{Glosario de Términos}

\begin{description}[style=nextline, leftmargin=1.5cm]
    \item[\textbf{Aplicación web}] Sistema de software accesible mediante un
    navegador web a través del protocolo HTTP/HTTPS, compuesto típicamente por
    un frontend, un backend y una base de datos.

    \item[\textbf{Ataque de inyección}] Técnica de explotación en la que un
    atacante inserta código malicioso en los datos de entrada de una aplicación
    para alterar su comportamiento o acceder a información no autorizada.

    \item[\textbf{Caja negra}] Metodología de pruebas en la que el evaluador
    no tiene acceso al código fuente ni a la arquitectura interna del sistema,
    simulando la perspectiva de un atacante externo.

    \item[\textbf{Crawler}] Componente de software que recorre automáticamente
    las páginas de un sitio web, descubriendo enlaces, formularios y puntos
    de entrada para su posterior análisis.

    \item[\textbf{Endpoint}] Punto de acceso específico de una aplicación web
    (URL) que acepta peticiones HTTP y retorna respuestas.

    \item[\textbf{Exploit}] Código o técnica que aprovecha una vulnerabilidad
    específica para comprometer la seguridad de un sistema.

    \item[\textbf{Falso positivo}] Resultado de una prueba de seguridad que
    indica incorrectamente la presencia de una vulnerabilidad que en realidad
    no existe.

    \item[\textbf{Fingerprinting}] Proceso de identificación de las tecnologías
    subyacentes de un servidor web (sistema operativo, servidor HTTP, lenguaje
    de programación, CMS) mediante el análisis de sus respuestas.

    \item[\textbf{Framework}] Estructura de software reutilizable que proporciona
    una arquitectura base sobre la cual se construyen aplicaciones específicas.

    \item[\textbf{Modelo de lenguaje (LLM)}] Sistema de inteligencia artificial
    entrenado sobre grandes volúmenes de texto que puede generar, analizar y
    contextualizar información en lenguaje natural.

    \item[\textbf{Módulo}] Componente de software independiente y autocontenido
    que implementa una funcionalidad específica dentro de un sistema más amplio.

    \item[\textbf{Payload}] Datos o código específico enviado a una aplicación
    web con el objetivo de provocar un comportamiento anómalo que revele una
    vulnerabilidad.

    \item[\textbf{Pentesting}] Prueba de penetración: evaluación de seguridad
    autorizada que simula ataques reales para identificar vulnerabilidades
    explotables.

    \item[\textbf{Superficie de ataque}] Conjunto total de puntos de entrada
    y vectores a través de los cuales un atacante puede intentar comprometer
    un sistema.

    \item[\textbf{Vulnerabilidad}] Debilidad o defecto en el diseño,
    implementación o configuración de un sistema que puede ser explotada para
    comprometer su seguridad.
\end{description}

% ===========================================================================
% CUERPO DEL TRABAJO
% ===========================================================================
\newpage
\pagenumbering{arabic}

% ===========================================================================
% CAPÍTULO I — EL PROBLEMA
% ===========================================================================
\section{CAPÍTULO I: EL PROBLEMA}

\subsection{Introducción}

En el panorama tecnológico actual, la integridad de las aplicaciones web es
un componente esencial para la operatividad de cualquier organización. Las
aplicaciones web se han convertido en el principal canal de interacción entre
las empresas y sus clientes, gestionando desde transacciones financieras
hasta datos personales sensibles. Sin embargo, esta ubicuidad digital ha
transformado a las aplicaciones web en el vector de ataque más explotado
por actores maliciosos a nivel global.

Las cifras recientes son contundentes. Según el informe Verizon Data Breach
Investigations Report (\cite{verizon2024}), la explotación de
vulnerabilidades en aplicaciones web experimentó un incremento del 180\%
durante 2023, con 1.997~incidentes clasificados como ataques a aplicaciones
web básicas, de los cuales 881 resultaron en divulgación confirmada de datos.
El informe IBM Cost of a Data Breach (\cite{ibm2024}) cifra el costo
promedio global de una brecha de datos en \$4,88~millones~USD, un aumento
del 10\% respecto al año anterior. Paralelamente, el informe Imperva Bad Bot
Report (\cite{imperva2025}) revela que por primera vez en la historia, el
tráfico automatizado de bots superó al tráfico humano, alcanzando el 51\%
del total global, con aproximadamente el 40\% clasificado como tráfico
malicioso.

Ante el incremento exponencial de ciberataques automatizados, surge la
necesidad imperiosa de contar con mecanismos de auditoría eficientes,
accesibles y adaptables. Esta urgencia se ha intensificado debido al auge
del desarrollo de software asistido por Inteligencia Artificial: estudios
recientes indican que entre el 45\% y el 48\% del código generado por
modelos de lenguaje contiene fallos de seguridad, con una densidad de
vulnerabilidades 2,7 veces superior al código escrito por humanos. Hasta
el 92\% de las aplicaciones construidas con asistencia de IA contienen al
menos una vulnerabilidad crítica, evidenciando que la aceleración del
desarrollo no va acompañada de las garantías de seguridad necesarias.

Bajo esta premisa, la presente investigación propone el diseño y
construcción de una herramienta de seguridad defensiva ---VulnSeeker---
ajustada a las realidades económicas y técnicas locales. A continuación, se
exponen los fundamentos teóricos previos, la relevancia social del proyecto
y la descripción detallada de la problemática que esta solución tecnológica
busca resolver.

\subsection{Planteamiento del Problema}

En el ecosistema digital actual, la seguridad de las aplicaciones web es un
requisito crítico. Sin embargo, existe una brecha de seguridad significativa
que afecta desproporcionadamente a desarrolladores independientes y Pequeñas y
Medianas Empresas (PYMEs). Estas entidades a menudo carecen del presupuesto
necesario para adquirir licencias de herramientas de auditoría comercial de alto
nivel (como Burp Suite Professional o Nessus), cuyos costos pueden ser
prohibitivos. Como consecuencia, muchas aplicaciones web son desplegadas en
entornos de producción sin una validación de seguridad adecuada, quedando
expuestas a vulnerabilidades críticas.

Desde una perspectiva arquitectónica, toda aplicación web puede concebirse como
una \textit{edificación digital} cuya integridad depende de tres propiedades
estructurales. En primer lugar, el \textbf{perímetro de seguridad}: el conjunto
de fronteras lógicas ---firewalls, reglas de acceso, políticas de
autenticación--- que delimitan el interior protegido del exterior hostil, de
manera análoga a los muros perimetrales de una construcción física. En segundo
lugar, la \textbf{superficie de exposición}: la totalidad de puntos de
interacción (formularios, APIs, endpoints públicos) que constituyen las
\textit{fachadas} por las cuales un atacante puede intentar el acceso; a mayor
superficie, mayor riesgo, del mismo modo en que un edificio con más ventanas
requiere más puntos de cierre. En tercer lugar, la \textbf{resiliencia
estructural}: la capacidad de la aplicación para mantener su operación crítica
aun cuando un componente individual sea comprometido, equivalente a la
redundancia estructural que impide que el fallo de una viga provoque el colapso
global. Cuando una PYME despliega una aplicación sin auditar su perímetro, sin
mapear su superficie de exposición y sin evaluar su resiliencia, está, en
términos arquitectónicos, habilitando una edificación sin inspección técnica.

Esta situación plantea la necesidad de una solución tecnológica accesible y
eficiente que cierre esta brecha de seguridad. Ante el incremento exponencial de
ciberataques automatizados, surge la necesidad imperiosa de contar con mecanismos
de auditoría eficientes, accesibles y adaptables, ajustados a las realidades
económicas y técnicas locales. VulnSeeker se propone como esa herramienta de
\textit{inspección técnica digital}: un instrumento que evalúa sistemáticamente
el perímetro, cuantifica la superficie de exposición y diagnostica la resiliencia
estructural de las aplicaciones web, democratizando un proceso que hasta ahora
estaba reservado a organizaciones con capacidad financiera para contratar
auditorías comerciales.

\subsection{Objetivo General}

Desarrollar un escáner modular para la detección automatizada de vulnerabilidades
en aplicaciones web.

\subsection{Objetivos Específicos}

\begin{enumerate}
    \item Analizar técnica y funcionalmente las herramientas de escaneo de
    vulnerabilidades de código abierto existentes.

    \item Diseñar la arquitectura de software modular de la herramienta,
    definiendo sus componentes principales.

    \item Implementar los módulos de detección para vulnerabilidades críticas
    iniciales como SQLi y XSS.

    \item Validar la efectividad de la herramienta desarrollada mediante la
    ejecución de pruebas de caja negra controladas.
\end{enumerate}

\subsection{Justificación}

La presente investigación se justifica en la necesidad de democratizar el acceso
a la ciberseguridad. En la actualidad, esta urgencia se ha intensificado debido
al auge del desarrollo de software asistido por Inteligencia Artificial (IA);
si bien estas herramientas aceleran drásticamente los tiempos de producción, a
menudo generan código que carece de validaciones y saneamiento riguroso, lo que
ha provocado una preocupante proliferación de aplicaciones web y sistemas
inherentemente inseguros.

Las cifras del panorama actual refuerzan esta justificación: el informe ENISA
Threat Landscape (\cite{enisa2025}) documenta que la explotación de
vulnerabilidades constituye el 21,3\% de los métodos de acceso inicial en
incidentes de ciberseguridad, mientras que los atacantes logran convertir nuevas
vulnerabilidades en exploits funcionales en cuestión de días tras su divulgación.
El costo promedio de una brecha de datos asciende a \$4,88~millones~USD según IBM
(\cite{ibm2024}), una cifra que para una PYME puede representar la diferencia
entre la continuidad operativa y el cese de actividades.

Frente a este panorama, la herramienta propuesta ofrece una solución de auditoría
automatizada, modular y de código abierto, eliminando las barreras económicas que
actualmente impiden a las PYMEs proteger sus activos digitales. Al reducir la
superficie de ataque en aplicaciones de pequeñas empresas, el proyecto contribuye
directamente a la seguridad del ecosistema digital local y nacional, fomentando
una cultura de desarrollo seguro. Adicionalmente, el contexto regional agrava la
urgencia: según el informe conjunto de la OEA, el BID y la Universidad de Oxford
(\cite{oasidb2025}), la mayoría de los países de América Latina y el Caribe
permanecen en un nivel de madurez en ciberseguridad medio o inferior, con brechas
estructurales en financiamiento, desarrollo de talento y coordinación
intersectorial.

\subsection{Viabilidad}

El desarrollo de la herramienta propuesta es técnica, operativa y económicamente
viable. En el aspecto técnico, se cuenta con el dominio de los lenguajes de
programación requeridos y el acceso a tecnologías de código abierto necesarias
para la construcción del escáner. Operativamente, se dispone de la asesoría
académica especializada del tutor en el área de Inteligencia Artificial y
Ciberseguridad, así como del equipo de computación necesario para los entornos
de prueba. Económicamente, el proyecto es factible al basarse en software libre,
eliminando costos de licenciamiento.

% ===========================================================================
% CAPÍTULO II — MARCO TEÓRICO
% ===========================================================================
\newpage
\section{CAPÍTULO II: MARCO TEÓRICO}

\subsection{Antecedentes de la Investigación}

Abdulghaffar et al. (\citeyear*{abdulghaffar2023}) publicaron un artículo de
revista científica titulado \textit{Enhancing Web Application Security through
Automated Penetration Testing with Multiple Vulnerability Scanners} en la
revista Computers (MDPI), desarrollado en colaboración entre la Glasgow
Caledonian University y la Birmingham City University (Reino Unido). Su objetivo
fue mejorar la detección de vulnerabilidades integrando varios motores. Bajo la
metodología de desarrollo de cascada iterativa, crearon un framework que
implementa algoritmos de combinación de resultados. Los hallazgos confirmaron
que la fusión de motores supera la detección individual, aunque eleva la
complejidad de las alertas. Su aporte valida técnicamente el algoritmo de
orquestación y la justificación tecnológica para el escáner modular propuesto
en este proyecto.

Así mismo, Alhogail y Alkahtani (\citeyear*{alhogail2024}) presentaron un
artículo de conferencia titulado \textit{Automated Extension-Based Penetration
Testing for Web Vulnerabilities}, publicado en las actas de la 15th ANT y
recopilado en Procedia Computer Science (King Saud University, Arabia Saudita).
El estudio buscó validar un modelo de pentesting automatizado basado en extensiones.
Emplearon una metodología experimental para prototipar Kashef bajo una arquitectura
de tres capas, validada mediante pruebas unitarias y funcionales. Los resultados
probaron que el enfoque extensible mejora la escalabilidad sin alterar el núcleo,
aunque carece de control de versiones. Este estudio fundamenta la arquitectura
modular y adaptable del presente proyecto, confirmando la viabilidad técnica de
incorporar nuevos módulos sin reestructurar el sistema base.

De igual modo, Arta et al. (\citeyear*{arta2023}) elaboraron un artículo de
investigación titulado \textit{Vulnerability Analysis and Effectiveness of OWASP
ZAP and Arachni on Web Security Systems} (Universitas Islam Riau, Indonesia),
comparando la eficacia de ambos escáneres. Mediante un análisis experimental
secuencial y usando OWASP Benchmark como estándar, midieron métricas de precisión
y cobertura en entornos controlados. Los resultados indicaron que ZAP alcanzó
tasas de detección superiores en la mayoría de escenarios. Este antecedente aporta
métricas empíricas que justifican la selección de herramientas de código abierto
como referencia para esta investigación.

Por otra parte, Repudi y Harinadh (\citeyear*{repudi2024}) difundieron un artículo
de revista científica titulado \textit{AutoPent: An Advanced Automated Web
Vulnerability Scanner} en el IRJEDT (Universal College of Engineering, India),
construyendo una herramienta automatizada de arquitectura modular. Utilizaron un
enfoque de desarrollo por capas (Interfaz, Lógica, Datos, Reportes) con
procesamiento multihilo. Los resultados validaron una reducción del 40\% en
tiempos de escaneo y una detección efectiva de vulnerabilidades comunes. Su aporte
constituye una referencia directa para la estructuración del flujo de escaneo por
capas y la generación de reportes, avalando la decisión de diseño modular y
concurrente planteada en la presente tesis.

A su vez, Savova et al. (\citeyear*{savova2024}) presentaron un artículo de
revista científica titulado \textit{Automated Web Application Scanning with
Wapiti, Selenium, and SQLMap} en la revista Security \& Future (National Military
University, Bulgaria), comparando la integración de herramientas open source.
Aplicaron una metodología de benchmarking según la Web Security Testing Guide
v4.2, evaluando cobertura y velocidad. Concluyeron que combinar Wapiti, Selenium
y SQLMap optimiza la auditoría compensando debilidades individuales. El estudio
guía el diseño arquitectónico de este proyecto al validar la orquestación de
funciones especializadas en componentes modulares definidos.

Igualmente, Yarphel y Rani (\citeyear*{yarphel2025}) publicaron un artículo de
revista científica titulado \textit{Comparative Performance Analysis of Web
Vulnerability Scanners} en el International Journal of Information Security
(Lovely Professional University, India), analizando el rendimiento de escáneres
comerciales y libres. Mediante investigación experimental cuantitativa y el Índice
de Youden, midieron la detección en escenarios locales y en vivo. Hallaron que ZAP
equilibra mejor precisión y usabilidad, aunque ninguna herramienta cubre todo. Este
antecedente reciente justifica empíricamente la necesidad de un sistema modular que
integre varios motores, demostrando con datos actuales que la unificación es
requerida para garantizar una cobertura integral.

Finalmente, Altulaihan et al. (\citeyear*{altulaihan2023}) realizaron un artículo
de revisión titulado \textit{A Survey on Web Application Penetration Testing} en
la revista Electronics (MDPI) (King Faisal University, Arabia Saudita), analizando
el estado del arte en pentesting web. Utilizaron la metodología de revisión
sistemática PRISMA, filtrando estudios bajo criterios del OWASP Top~10.
Demostraron que las herramientas actuales suelen carecer de modularidad o cobertura
total frente a fallos críticos. Su aporte fundamenta teóricamente los vacíos
técnicos existentes, legitimando la pertinencia científica de desarrollar una nueva
herramienta modular, flexible y adaptada a entornos tecnológicos actuales como la
propuesta en esta investigación.

\subsection{Panorama Global de Amenazas}

La adopción masiva de arquitecturas web para soportar la operatividad de
organizaciones de toda escala ha venido acompañada de un incremento correlativo
en la frecuencia y sofisticación de los incidentes de seguridad digital. El
estudio sistemático de los reportes anuales de las firmas y agencias de
ciberseguridad de mayor relevancia internacional permite constatar que las
aplicaciones web constituyen en la actualidad el principal vector de intrusión a
nivel global.

De acuerdo con el informe \textit{Verizon Data Breach Investigations Report}
(\cite{verizon2024}), el compromiso de aplicaciones web experimentó un
crecimiento del 180\% durante el año 2023, consolidándose como la categoría
predominante de ataque dentro de los incidentes de seguridad básicos. Este
reporte documenta 1.997~incidentes específicos dirigidos a aplicaciones web,
de los cuales 881 resultaron en una divulgación confirmada de datos sensibles.
Este incremento se atribuye en gran medida a la persistencia de debilidades en
la capa de presentación e interfaces de programación de aplicaciones (APIs), las
cuales son explotadas sistemáticamente para la exfiltración de bases de datos
y ejecución remota de comandos.

Desde el punto de vista del impacto financiero, el informe \textit{Cost of a Data
Breach Report} de IBM Security (\cite{ibm2024}) establece que el costo promedio
global de una brecha de datos alcanzó los \$4,88~millones~USD en
2024, lo que representa un incremento del 10\% respecto al valor reportado en
2023. Este incremento marca la mayor tasa de crecimiento interanual de los
costos de incidentes desde la pandemia del COVID-19. Para el segmento de las
Pequeñas y Medianas Empresas (PYMEs), que operan con márgenes financieros
estrechos, un incidente de esta magnitud resulta destructivo, comprometiendo no
solo su flujo de caja debido a penalizaciones y costos de remediación, sino
también la confianza reputacional necesaria para su continuidad operativa.

Adicionalmente, el análisis del tráfico automatizado en la red aporta datos sobre
la masificación del riesgo. Según el \textit{Bad Bot Report} publicado por
Imperva (\cite{imperva2025}), el tráfico originado por agentes automatizados
superó por primera vez en la historia al tráfico legítimo generado por humanos,
alcanzando el 51\% del volumen global. Dentro de esta porción de tráfico,
aproximadamente el 40\% es clasificado como tráfico malicioso de bots, el cual se
dedica a la exploración masiva de puertos, rastreo de directorios públicos
y ataques de fuerza bruta. Esta automatización constante implica que cualquier
aplicación expuesta en internet, independientemente de su relevancia comercial,
es escaneada e intentada vulnerar en cuestión de horas tras su puesta en línea,
invalidando la suposición común de que las PYMEs no constituyen objetivos
relevantes para los ciberatacantes.

Finalmente, la Agencia de la Unión Europea para la Ciberseguridad (ENISA), en su
informe \textit{Threat Landscape 2025} (\cite{enisa2025}), destaca que la
explotación de vulnerabilidades de software conocidas es responsable del 21,3\%
de los accesos iniciales exitosos en intrusiones de red. La velocidad con la
que los actores de amenazas desarrollan exploits funcionales tras la divulgación de
una vulnerabilidad se ha reducido a un promedio de tres días, lo que exige una
capacidad de auditoría defensiva constante y ágil que permita identificar
debilidades antes de que sean explotadas comercialmente.

% ── Bases teóricas ──────────────────────────────────────────────────────
\subsection{Bases Teóricas}


\subsubsection{Seguridad en Aplicaciones Web}

La seguridad en aplicaciones web comprende el conjunto de prácticas, técnicas y
herramientas orientadas a proteger los sistemas web contra amenazas y
vulnerabilidades que puedan comprometer la confidencialidad, integridad y
disponibilidad de la información. Según la OWASP Foundation
(\citeyear*{owasp2021}), las aplicaciones web representan el vector de ataque
más explotado en la actualidad, debido a su exposición directa en internet y a
la complejidad inherente de sus arquitecturas multicapa.

La superficie de ataque de una aplicación web incluye todos los puntos de
interacción entre el usuario y el servidor: formularios de entrada, parámetros
de URL, cabeceras HTTP, cookies de sesión, APIs REST y mecanismos de
autenticación. Un modelo de amenazas efectivo debe considerar tanto los ataques
automatizados (bots, escáneres) como los ataques dirigidos por actores
sofisticados.

\subsubsection{Pruebas de Seguridad de Aplicaciones Dinámicas (DAST)}

Las Pruebas de Seguridad de Aplicaciones Dinámicas (DAST, por sus siglas en
inglés) constituyen una metodología de evaluación de caja negra que analiza las
aplicaciones web en tiempo de ejecución, sin acceso al código fuente. A
diferencia de las Pruebas de Seguridad de Aplicaciones Estáticas (SAST), que
inspeccionan el código sin ejecutarlo, y de las Pruebas Interactivas (IAST),
que instrumentan el runtime, el enfoque DAST simula el comportamiento de un
atacante externo real.

\begin{table}[H]
    \centering
    \caption{Comparación de metodologías de testing de seguridad}
    \label{tab:sast-dast-iast}
    \begin{tabular}{p{3cm} p{3.5cm} p{3.5cm} p{3.5cm}}
        \toprule
        \textbf{Criterio} & \textbf{SAST} & \textbf{DAST} & \textbf{IAST} \\
        \midrule
        Acceso al código & Sí & No & Sí (instrumentado) \\
        Fase de ejecución & Pre-compilación & Runtime & Runtime \\
        Falsos positivos & Alto & Bajo & Bajo \\
        Cobertura & Código analizado & Endpoints expuestos & Ambos \\
        Ejemplo & SonarQube & \textbf{VulnSeeker}, ZAP & Contrast \\
        \bottomrule
    \end{tabular}
\end{table}

VulnSeeker opera exclusivamente bajo el paradigma DAST, enviando peticiones HTTP
reales al objetivo y analizando las respuestas para identificar patrones
indicativos de vulnerabilidades.

\subsubsection{OWASP Top 10 (2021)}

El OWASP Top~10 es un documento de referencia publicado por la Open Worldwide
Application Security Project (OWASP) que cataloga las diez categorías de riesgos
de seguridad más críticos en aplicaciones web. La edición de 2021
(\cite{owasp2021}) reorganizó las categorías con base en datos de incidentes
reales recopilados de más de 500.000 aplicaciones. Las diez categorías son:

\begin{enumerate}
    \item \textbf{A01 — Broken Access Control:} Fallos en la restricción de
    acceso a recursos protegidos. VulnSeeker lo evalúa mediante los módulos
    Open Redirect y CORS Scanner.

    \item \textbf{A02 — Cryptographic Failures:} Exposición de datos sensibles
    por criptografía débil o ausente. Cubierto por TLS Checker y Sensitive Data
    Scanner.

    \item \textbf{A03 — Injection:} Inyección de código malicioso en intérpretes
    backend. VulnSeeker implementa cinco módulos: SQLi, XSS, Command Injection,
    LFI y RFI.

    \item \textbf{A04 — Insecure Design:} Deficiencias en el diseño de la
    aplicación. Evaluado parcialmente por File Upload Detector.

    \item \textbf{A05 — Security Misconfiguration:} Configuraciones inseguras
    del servidor. Cubierto por Header Analyzer, Dir Listing, HTTP Methods
    y Path Fuzzer.

    \item \textbf{A06 — Vulnerable and Outdated Components:} Uso de componentes
    con vulnerabilidades conocidas. VulnSeeker consulta la NVD API del NIST en
    tiempo real mediante CVE Lookup.

    \item \textbf{A07 — Identification and Authentication Failures:} Fallos en
    autenticación y gestión de sesiones. Cubierto por CSRF Auditor, Brute Force
    Detector y Weak Session Auditor.

    \item \textbf{A08 — Software and Data Integrity Failures:} Fallos en la
    integridad de datos y actualizaciones. No evaluable externamente mediante
    DAST.

    \item \textbf{A09 — Security Logging and Monitoring Failures:} Deficiencias
    en registro y monitoreo. No evaluable externamente mediante DAST.

    \item \textbf{A10 — Server-Side Request Forgery (SSRF):} Manipulación del
    servidor para realizar solicitudes a recursos internos. Cubierto por el
    módulo SSRF Scanner.
\end{enumerate}

VulnSeeker cubre 8 de las 10 categorías. Las categorías A08 y A09 requieren
acceso a la infraestructura interna del servidor, lo cual escapa al alcance de
una herramienta DAST de caja negra.

\subsubsection{Evaluación de Efectividad y Metodología OWASP Benchmark}

La validación científica de una herramienta de escaneo dinámico de seguridad
(DAST) exige el uso de marcos de referencia empíricos que permitan cuantificar
su precisión de forma estandarizada. El estándar más extendido en la
industria es el \textit{OWASP Benchmark Project} (\cite{owaspbenchmark}), un
proyecto de código abierto diseñado para evaluar la tasa de aciertos y errores de
cualquier analizador de vulnerabilidades automatizado mediante una suite de miles
de casos de prueba controlados.

La efectividad de un escáner se mide a partir de la clasificación de sus
hallazgos en cuatro cuadrantes de una matriz de confusión: Verdaderos Positivos
(VP), Falsos Positivos (FP), Verdaderos Negativos (VN) y Falsos Negativos (FN).
A partir de estos valores se derivan tres métricas fundamentales de
rendimiento estadístico:

\begin{enumerate}
    \item \textbf{Tasa de Verdaderos Positivos (TPR - True Positive Rate):}
    Denominada también como sensibilidad o recall, mide la proporción de
    vulnerabilidades reales que fueron detectadas exitosamente por el escáner:
    \[ TPR = \frac{VP}{VP + FN} \]

    \item \textbf{Tasa de Falsos Positivos (FPR - False Positive Rate):} Mide
    la proporción de casos limpios o no vulnerables que fueron erróneamente
    catalogados como vulnerabilidades por la herramienta:
    \[ FPR = \frac{FP}{FP + VN} \]

    \item \textbf{Índice de Youden ($J$):} Constituye la métrica integradora
    del OWASP Benchmark. Evalúa la capacidad neta de discriminación de la
    herramienta, restando la tasa de error a la tasa de acierto:
    \[ J = TPR - FPR \]
    El valor de $J$ oscila entre $-1$ y $+1$. Un valor de $+1$ representa un
    escáner perfecto (100\% de verdaderos positivos y 0\% de falsos positivos),
    mientras que un valor de $0$ o menor indica que el comportamiento de la
    herramienta es estadísticamente equivalente o inferior al de una selección
    completamente aleatoria.
\end{enumerate}

El OWASP Benchmark demuestra empíricamente que muchos escáneres DAST
comerciales y de código abierto sufren de tasas elevadas de falsos positivos,
lo que genera fatiga de alertas en los equipos de desarrollo. Este marco
conceptual orienta el diseño de VulnSeeker, que prioriza la minimización de la
tasa de falsos positivos ($FPR \to 0$) en sus módulos de inyección activa
mediante reglas de validación multi-paso, comparación contra respuestas
\textit{baseline} y análisis de firmas de error, incluso si ello implica una
reducción menor en la tasa general de cobertura. La aplicación formal del OWASP
Benchmark para obtener el Índice de Youden cuantitativo del framework se plantea
como línea de validación futura (véase la sección de Recomendaciones), dado que
excede el alcance experimental del presente trabajo, centrado en entornos de
práctica estándar.

\subsubsection{Arquitectura Modular de Software}


La arquitectura modular es un enfoque de diseño que descompone un sistema en
componentes independientes y cohesivos, cada uno con una responsabilidad bien
definida. Bass, Clements y Kazman (\citeyear*{bass2012}) definen la
arquitectura de software como el vehículo principal para alcanzar los
atributos de calidad de un sistema, argumentando que la funcionalidad y los
atributos de calidad son ortogonales: mientras la funcionalidad define
\textit{qué} hace el sistema, la arquitectura determina \textit{qué tan bien}
lo hace en términos de modificabilidad, rendimiento, disponibilidad,
seguridad y testabilidad.

Este paradigma se fundamenta en los principios SOLID formulados por
Martin (\citeyear*{martin2003}) y posteriormente elevados a escala
arquitectónica en su obra sobre arquitectura limpia (\cite{martin2017}),
donde establece la \textit{Regla de Dependencia}: las dependencias del
código fuente deben apuntar únicamente hacia adentro, hacia las políticas
de más alto nivel, garantizando que las reglas de negocio permanezcan
independientes de la interfaz de usuario, las bases de datos y los
frameworks externos:

\begin{itemize}
    \item \textbf{Single Responsibility:} Cada módulo tiene una única razón de
    cambio. En VulnSeeker, cada módulo detecta un tipo específico de
    vulnerabilidad.

    \item \textbf{Open/Closed:} El sistema es abierto a la extensión pero
    cerrado a la modificación. Nuevos módulos se agregan sin alterar el motor
    central.

    \item \textbf{Liskov Substitution:} Todo módulo que herede de
    \texttt{ScannerModule} es intercambiable en el motor de ejecución.

    \item \textbf{Interface Segregation:} La interfaz base define solo los métodos
    necesarios: \texttt{scan()} y metadatos del módulo.

    \item \textbf{Dependency Inversion:} El motor depende de la abstracción
    \texttt{ScannerModule}, no de implementaciones concretas.
\end{itemize}

VulnSeeker implementa el patrón Strategy tal como lo definen Gamma et~al.
(\citeyear*{gamma1995}): \textit{``Definir una familia de algoritmos,
encapsular cada uno, y hacerlos intercambiables. Strategy permite que el
algoritmo varíe independientemente de los clientes que lo utilizan''}
(p.~315). En VulnSeeker, el objeto Engine delega la ejecución a cada módulo
registrado sin conocer su lógica interna, garantizando alta cohesión y bajo
acoplamiento. Cada módulo de detección constituye una \textit{estrategia}
concreta que implementa la interfaz \texttt{ScannerModule}, permitiendo que
el motor de orquestación los trate de forma uniforme sin necesidad de
condicionales anidados ni conocimiento de la implementación subyacente.

Esta separación de responsabilidades se alinea también con los patrones
de Fowler (\citeyear*{fowler2002}), quien propone el desacoplamiento de
capas como estrategia fundamental para crear sistemas mantenibles, testables
y escalables: la capa de presentación (GUI), la capa de dominio (módulos de
detección) y la capa de acceso a datos (SQLite) operan de forma independiente,
comunicándose exclusivamente a través de interfaces bien definidas.

\subsubsection{Inteligencia Artificial y la Proliferación de Código Inseguro}

Un fenómeno emergente que agrava la necesidad de herramientas DAST accesibles
es la proliferación de código inseguro generado por modelos de lenguaje de
gran escala. Investigaciones recientes indican que entre el 45\% y el 48\%
del código producido por asistentes de programación basados en IA contiene
fallos de seguridad, con una densidad de vulnerabilidades 2,7~veces superior
al código escrito por humanos. Las deficiencias más frecuentes incluyen la
ausencia de validación de entradas (\textit{input sanitization}), secretos
hardcodeados en el código fuente, fallos criptográficos y vulnerabilidades
de inyección de logs.

Esta realidad tiene implicaciones directas para el ecosistema de las PYMEs:
muchos desarrolladores independientes y equipos pequeños han adoptado
herramientas de generación de código asistido por IA para acelerar sus
ciclos de desarrollo, produciendo aplicaciones que funcionalmente cumplen
sus requisitos pero que estructuralmente carecen de las validaciones de
seguridad necesarias. El resultado es una proliferación de aplicaciones
web desplegadas en producción con vulnerabilidades conocidas que podrían
ser detectadas y remediadas mediante un escaneo DAST automatizado como el
que VulnSeeker provee.

\subsubsection{Contexto Regional: Ciberseguridad en América Latina}

El contexto regional agrega una dimensión adicional de urgencia a la
presente investigación. Según el informe conjunto de la Organización de
Estados Americanos (OEA), el Banco Interamericano de Desarrollo (BID) y
la Universidad de Oxford (\cite{oasidb2025}), la región de América Latina
y el Caribe presenta una trayectoria positiva en madurez de ciberseguridad,
pero la mayoría de los países permanecen en un nivel medio o inferior en
el Modelo de Madurez de Ciberseguridad (CMM). Las brechas más significativas
se concentran en financiamiento, desarrollo de talento especializado y
coordinación intersectorial.

En el caso específico de Venezuela, el panorama presenta particularidades
que refuerzan la pertinencia de VulnSeeker. En agosto de 2024, el gobierno
venezolano estableció el Consejo Nacional de Ciberseguridad mediante el
Decreto N°~4.975 (\cite{venezuelacyber2024}), con la misión de desarrollar
políticas, promover la capacitación y establecer una red de vigilancia para
incidentes telemáticos. Sin embargo, el país aún carece de una estrategia
nacional de ciberseguridad plenamente formalizada e integrada, lo que
evidencia un vacío institucional que herramientas de código abierto como
VulnSeeker pueden contribuir a mitigar desde la perspectiva técnica.

Las principales amenazas identificadas en la región incluyen el auge del
ransomware, el robo de datos y los ataques dirigidos a infraestructura
crítica (energía, telecomunicaciones). En este contexto, la disponibilidad
de una herramienta DAST gratuita, modular y adaptable a las realidades
técnicas locales constituye una contribución directa a la resiliencia
digital del tejido empresarial latinoamericano.

\subsubsection{Common Vulnerabilities and Exposures (CVE) y Taxonomías Asociadas}

El ecosistema de identificación de vulnerabilidades se sustenta en tres
taxonomías complementarias mantenidas por organismos internacionales:

\begin{itemize}
    \item \textbf{CVE (Common Vulnerabilities and Exposures):} Catálogo
    estandarizado mantenido por MITRE Corporation que asigna un identificador
    único (ej: CVE-2021-44228) a cada vulnerabilidad divulgada públicamente.

    \item \textbf{CWE (Common Weakness Enumeration):} Taxonomía de debilidades
    de software que categoriza los tipos de errores de programación (ej:
    CWE-89 para SQL Injection).

    \item \textbf{CVSS (Common Vulnerability Scoring System):} Sistema de
    puntuación estandarizado que asigna una severidad numérica (0.0 a 10.0) a
    cada vulnerabilidad, clasificándola como None, Low, Medium, High o Critical.
\end{itemize}

VulnSeeker integra estas taxonomías a través de su módulo CVE Lookup, que
consulta la API 2.0 de la National Vulnerability Database (NVD) del NIST en
tiempo real. Cuando el módulo de fingerprinting identifica una tecnología y su
versión, CVE Lookup busca automáticamente las vulnerabilidades conocidas
asociadas, complementando la detección dinámica con inteligencia de amenazas
actualizada.

\subsubsection{Inteligencia Artificial Aplicada a la Ciberseguridad}

La aplicación de modelos de lenguaje de gran escala (LLMs) en ciberseguridad
representa una frontera emergente en la automatización de la auditoría de
seguridad. Los LLMs, entrenados sobre vastos corpus de documentación técnica,
reportes de incidentes y bases de conocimiento de vulnerabilidades, poseen la
capacidad de contextualizar hallazgos técnicos y traducirlos a evaluaciones de
riesgo comprensibles para audiencias no técnicas.

VulnSeeker integra el modelo Llama 3.3 70B de Meta, accedido a través de la API
de Groq, para generar informes ejecutivos que incluyen: nivel de riesgo global,
las cinco amenazas más críticas con explicación contextualizada, un plan de
acción inmediato (48 horas) y recomendaciones estratégicas a mediano plazo. Este
enfoque integra análisis asistido por inteligencia artificial generativa dentro
de una arquitectura DAST modular.


% ── Análisis comparativo de herramientas ────────────────────────────────
\subsection{Análisis Comparativo de Herramientas Existentes}

Para fundamentar las decisiones de diseño de VulnSeeker, se realizó un análisis
comparativo de las herramientas DAST más representativas del ecosistema actual.
La selección incluyó dos herramientas de código abierto ampliamente adoptadas
(OWASP ZAP y Nikto), una herramienta basada en plantillas (Nuclei) y una
herramienta comercial líder (Burp Suite Professional), con el fin de cubrir el
espectro completo de soluciones disponibles.

\begin{table}[H]
    \centering
    \caption{Matriz comparativa de herramientas DAST}
    \label{tab:comparativa}
    \small
    \begin{tabular}{p{2.5cm} p{2cm} p{2cm} p{2cm} p{2cm} p{2.2cm}}
        \toprule
        \textbf{Criterio} & \textbf{ZAP} & \textbf{Nikto} & \textbf{Nuclei} &
        \textbf{Burp Suite} & \textbf{VulnSeeker} \\
        \midrule
        Arquitectura &
            Plugin-based & Monolítica & Templates & Monolítica &
            Modular \\
        Extensibilidad &
            Alta (Java) & Baja & Alta (YAML) & Media (BApps) &
            Alta (Python) \\
        Costo anual &
            Gratis & Gratis & Gratis & \$449 USD &
            Gratis \\
        Análisis IA &
            No & No & No & No &
            Sí (Llama 3.3) \\
        CVE en vivo &
            No & No & Parcial & No &
            Sí (NVD API) \\
        Interfaz &
            GUI (Java) & CLI & CLI & GUI (Java) &
            GUI (Python) \\
        Crawler &
            Sí & No & No & Sí &
            Sí (con auth) \\
        Reportes PDF &
            Sí & No & No & Sí &
            Sí \\
        \bottomrule
    \end{tabular}
\end{table}

Del análisis se concluye que VulnSeeker ofrece tres diferenciadores clave
respecto a las herramientas evaluadas: (1) integración nativa de análisis
asistido por IA para la generación de informes ejecutivos, (2) consulta en
tiempo real a la NVD del NIST para correlación de CVEs, y (3) una arquitectura
modular en Python que reduce la barrera de entrada para la extensión de la
herramienta, a diferencia de ZAP que requiere conocimiento de Java o Nuclei que
se limita a plantillas YAML declarativas.

\subsubsection{Lecciones Aprendidas de los Referentes}

Más allá de la comparación cuantitativa, el estudio de cada herramienta reveló
patrones de diseño ---tanto virtuosos como deficientes--- que informaron
directamente las decisiones arquitectónicas de VulnSeeker.

\textbf{OWASP ZAP} constituye, sin duda, la referencia más completa del
ecosistema open source. No obstante, su principal vicio de diseño reside en la
\textit{sobrecarga cognitiva} de su interfaz: más de 200 opciones de
configuración distribuidas en menús anidados, paneles simultáneos y una curva de
aprendizaje que desalienta al usuario no especializado. VulnSeeker resuelve este
problema mediante una interfaz de cuatro pestañas (Escaneo, Resultados,
Historial, Configuración) que reduce la fricción operativa sin sacrificar
funcionalidad.

\textbf{Nikto}, aunque eficiente en la detección de configuraciones inseguras,
adolece de un vicio estructural crítico: su arquitectura monolítica impide la
extensibilidad. Agregar un nuevo tipo de detección requiere modificar el código
fuente central, violando el principio Open/Closed. VulnSeeker supera esta
limitación mediante la abstracción \texttt{ScannerModule}, que permite
incorporar nuevos módulos sin alterar una sola línea del motor.

\textbf{Nuclei} introdujo una innovación valiosa: la definición de reglas de
detección mediante plantillas YAML, democratizando la creación de checks de
seguridad. Sin embargo, esta misma decisión impone un \textit{techo de
expresividad}: las lógicas de detección complejas ---como el análisis
multi-paso de sesiones débiles o la correlación de CVEs en vivo--- son
inviables en un lenguaje declarativo. VulnSeeker opta por módulos en Python
puro, otorgando poder expresivo completo a cada detector.

\textbf{Burp Suite Professional} representa el estándar de oro comercial, pero
su modelo de licenciamiento (\$449~USD/año) lo posiciona como una herramienta
inaccesible para el segmento que VulnSeeker busca atender. Además, su sistema
de extensiones (BApps) opera dentro de un sandbox con restricciones de API que
limitan la profundidad de las integraciones.

En síntesis, VulnSeeker no pretende reemplazar a estas herramientas, sino
resolver sus vicios de diseño específicos: la complejidad de ZAP, la rigidez de
Nikto, el techo expresivo de Nuclei y la inaccesibilidad económica de Burp
Suite. Cada decisión arquitectónica del framework es una respuesta directa a una
lección identificada en el análisis de referentes.

\subsubsection{Filosofías de Diseño Contrastadas}

El análisis de referentes revela que cada herramienta encarna una
\textit{filosofía de diseño} distinta, y comprender estas filosofías es
esencial para justificar el paradigma que VulnSeeker propone.

OWASP ZAP opera bajo una filosofía de \textit{completitud maximalista}: busca
ser la única herramienta que el auditor necesite, integrando proxy, spider,
fuzzer, scanner activo y pasivo en un solo ejecutable. Esta ambición, aunque
admirable, genera lo que en arquitectura se denomina un \textit{programa
funcional sobredimensionado}: un sistema que intenta resolver todas las
necesidades imaginables y, al hacerlo, se vuelve inaccesible para el usuario
que solo necesita una auditoría básica.

Burp Suite, por su parte, adopta la filosofía del \textit{ecosistema
cerrado}: una plataforma central propietaria que permite extensiones
controladas (BApps) dentro de los límites que el fabricante define. Es el
equivalente digital del modelo de franquicia arquitectónica: el diseño base
es impecable, pero la personalización está acotada por las reglas del
concesionario, y el acceso exige un tributo anual.

VulnSeeker propone una tercera vía: la filosofía del \textit{framework
liviano y modular}. En lugar de aspirar a la completitud de ZAP o al pulido
comercial de Burp, VulnSeeker se diseña como un \textit{chásis abierto}
sobre el cual se montan módulos especializados. Esta filosofía prioriza tres
valores: (1) la \textit{accesibilidad cognitiva}, reduciendo la interfaz a
lo esencial; (2) la \textit{extensibilidad soberana}, donde cualquier
desarrollador puede crear un módulo sin depender de un marketplace ni de una
licencia; y (3) el \textit{costo cero operativo}, garantizando que la única barrera de entrada sea el conocimiento técnico, no el capital financiero.
Este paradigma resulta particularmente idóneo para el ecosistema de las PYMEs
latinoamericanas, donde la restricción presupuestaria no debe implicar una
renuncia a la seguridad.

% ── Bases legales ───────────────────────────────────────────────────────
\subsection{Bases Legales}

El desarrollo y uso de herramientas de auditoría de seguridad informática se
enmarca en un contexto legal que regula tanto la protección de sistemas como la
legalidad de las pruebas de intrusión automatizadas.

\subsubsection{Marco Legal Venezolano}

En Venezuela, la \textbf{Ley Especial contra los Delitos Informáticos} (Gaceta
Oficial N° 37.313, 2001) tipifica como delito el acceso indebido a sistemas de
información (Art.~6), el sabotaje informático (Art.~7) y la violación de la
privacidad de datos (Art.~20). Sin embargo, la ley establece como elemento
constitutivo del delito la \textit{intención dolosa} y la \textit{ausencia de
autorización}. Las pruebas de seguridad realizadas con autorización explícita y
por escrito del propietario del sistema quedan fuera del ámbito penal, lo cual
legitima el uso de VulnSeeker en entornos controlados y con consentimiento
formal.

\subsubsection{Estándares Internacionales}

Los principales estándares que regulan y guían las auditorías de seguridad
informática incluyen:

\begin{itemize}
    \item \textbf{ISO/IEC 27001:2022:} Estándar internacional para la gestión
    de la seguridad de la información, que establece un marco de controles para
    identificar, evaluar y tratar riesgos de seguridad (\cite{iso27001}).

    \item \textbf{OWASP Testing Guide v4.2:} Guía metodológica publicada por la
    OWASP Foundation que define procedimientos detallados para evaluar la
    seguridad de aplicaciones web, incluyendo pruebas de inyección, autenticación
    y gestión de sesiones (\cite{owasptesting}).

    \item \textbf{PTES (Penetration Testing Execution Standard):} Estándar
    técnico que define las fases de una prueba de intrusión ética: definición del
    alcance, reconocimiento, análisis de vulnerabilidades, explotación y reporte
    (\cite{ptes}).
\end{itemize}

VulnSeeker opera bajo el principio del \textit{consentimiento informado}: la
herramienta incluye un disclaimer legal visible que advierte al operador sobre
la obligatoriedad de contar con autorización escrita antes de ejecutar cualquier
escaneo. El sistema no incorpora funcionalidades de explotación activa; su
alcance se limita a la detección y reporte de vulnerabilidades.

% ===========================================================================
% CAPÍTULO III — MARCO METODOLÓGICO
% ===========================================================================
\newpage
\section{CAPÍTULO III: MARCO METODOLÓGICO}

\subsection{Tipo y Diseño de Investigación}

La investigación se desarrolla bajo la modalidad de Proyecto Factible con un
diseño de campo experimental y documental. El sistema opera bajo un modelo de
Pruebas de Seguridad de Aplicaciones Dinámicas (DAST) con enfoque de Caja Negra.

En cuanto a la ingeniería de software, se adoptó el Modelo de Desarrollo
Evolutivo por Prototipos, fundamentado por Sommerville (\citeyear*{sommerville2011}),
quien establece que el prototipado evolutivo es un enfoque donde ``el objetivo es
trabajar con el cliente para explorar sus requisitos y entregar un sistema final.
El desarrollo comienza con las partes del sistema que están bien comprendidas. El
sistema se refina e incrementa a través de varias etapas basándose en la
retroalimentación del cliente hasta alcanzar el sistema final'' (p.~43). A
diferencia del prototipado desechable, el prototipo evolutivo está destinado a
evolucionar hacia el sistema de producción final.

Esta metodología resultó idónea para VulnSeeker por tres razones fundamentales:
(1)~la naturaleza incremental del desarrollo ---donde cada módulo de detección
puede ser implementado, probado y refinado de forma independiente antes de
integrarse al motor central---; (2)~la posibilidad de validación temprana
mediante escaneos parciales que retroalimentan el diseño de módulos posteriores;
y (3)~la compatibilidad con el patrón Strategy implementado, donde cada nuevo
módulo-prototipo extiende el sistema sin requerir modificaciones al núcleo
existente.

\subsection{Técnicas e Instrumentos de Recolección de Datos}

Las técnicas utilizadas para la recolección de datos en esta investigación
fueron seleccionadas en correspondencia con los objetivos específicos y la
naturaleza experimental del proyecto:

\begin{enumerate}
    \item \textbf{Revisión documental:} Análisis sistemático de la literatura
    académica, documentación técnica de herramientas existentes (OWASP ZAP,
    Nikto, Nuclei, Burp Suite) y estándares internacionales (OWASP Top~10,
    PTES, ISO~27001). Esta técnica fundamentó el diagnóstico del Objetivo~1
    y la construcción del marco teórico.

    \item \textbf{Benchmarking técnico:} Evaluación comparativa de las
    capacidades funcionales de las herramientas DAST del mercado mediante
    pruebas directas, utilizando como marco de referencia conceptual las
    categorías del OWASP Top~10 y las métricas definidas por el OWASP Benchmark
    Project (\cite{owaspbenchmark}): True Positive Rate (TPR), False
    Positive Rate (FPR) e Índice de Youden ($J = TPR - FPR$). La aplicación
    formal de este estándar como medición cuantitativa de VulnSeeker se reserva
    como validación futura.

    \item \textbf{Observación directa (escaneos controlados):} Ejecución
    de auditorías automatizadas contra entornos de práctica estándar (DVWA)
    con registro sistemático de hallazgos, tiempos de ejecución y métricas
    de detección.

    \item \textbf{Análisis de logs y métricas automatizadas:} Recolección
    y procesamiento de los registros generados por el framework durante los
    escaneos, incluyendo conteo de hallazgos por severidad, distribución por
    categoría OWASP, tiempos de respuesta y tasas de detección por módulo.

    \item \textbf{Pruebas unitarias automatizadas:} Validación del
    comportamiento de cada módulo mediante el framework \texttt{pytest},
    utilizando respuestas HTTP simuladas (mocking) para garantizar ejecución
    determinista e independiente de la red.
\end{enumerate}

\subsection{Población y Muestra}

La población usuaria de VulnSeeker comprende cuatro perfiles diferenciados,
cuyas necesidades técnicas y contextos operativos fundamentan las decisiones
de diseño de la herramienta.

\begin{table}[H]
    \centering
    \caption{Perfiles de usuario objetivo de VulnSeeker}
    \label{tab:perfiles}
    \small
    \begin{tabular}{p{2.5cm} p{4cm} p{3cm} p{3.5cm}}
        \toprule
        \textbf{Perfil} & \textbf{Necesidad técnica} &
        \textbf{Herramienta actual} & \textbf{Impacto de la brecha} \\
        \midrule
        Desarrollador independiente &
            Validar seguridad antes del deploy sin costo adicional &
            Ninguna / escaneos manuales &
            Ansiedad por vulnerabilidades, deuda técnica acumulada \\
        Auditor de seguridad junior &
            Herramienta complementaria automatizable para triaje inicial &
            ZAP (complejidad alta) &
            Ineficiencia por auditoría manual \\
        PYME ($<$~50 empleados) &
            Cumplir requisitos mínimos de seguridad sin equipo dedicado &
            Ninguna / consultores externos &
            Riesgo financiero y reputacional ante brechas \\
        Estudiante de computación &
            Aprendizaje práctico de técnicas de ciberseguridad &
            Herramientas dispersas &
            Brecha entre teoría académica y práctica profesional \\
        \bottomrule
    \end{tabular}
\end{table}

La muestra técnica para la validación experimental está compuesta por entornos
de práctica estándar de la industria: Damn Vulnerable Web Application (DVWA),
configurado con nivel de seguridad bajo para maximizar la superficie de ataque
detectable y permitir la evaluación completa de los 25~módulos.

\subsection{Criterios de Diseño del Framework}

El diseño de VulnSeeker se fundamenta en tres criterios rectores que,
análogamente a los principios de una obra arquitectónica, determinan la
idoneidad de la solución para su contexto de uso.

\subsubsection{Criterio Funcional: Eficiencia en la Detección}

Del mismo modo en que la función de un edificio define su distribución espacial,
la función primaria de VulnSeeker ---la detección de vulnerabilidades--- define
su flujo de datos. El criterio funcional exige que cada módulo de detección
opere con la máxima eficiencia: payloads específicos por tipo de inyección,
análisis de respuesta por patrones de error conocidos, y clasificación automática
por severidad siguiendo el estándar CVSS. Este criterio ga\-ran\-ti\-za que cada
interacción con el objetivo tenga un propósito diagnóstico preciso, evitando el
ruido generado por escaneos indiscriminados.

\subsubsection{Criterio Estructural: La Modularidad como Sistema de Soporte}

Así como la estructura portante de un edificio debe soportar cargas sin
transmitir esfuerzos indebidos a los elementos no estructurales, la arquitectura
modular de VulnSeeker garantiza que la adición o eliminación de un módulo no
genere esfuerzos de refactorización en el núcleo. La clase abstracta
\texttt{ScannerModule} actúa como la \textit{cimentación} del sistema: define
la interfaz contractual (\texttt{scan()}, metadatos, severidad) sobre la cual
se erigen los 25~módulos sin acoplamiento lateral. El motor de orquestación,
análogo a un sistema de vigas, distribuye la carga de ejecución entre hilos
concurrentes sin que ningún módulo conozca la existencia de los demás.

\subsubsection{Criterio de Usuario (UX): Simplicidad Operativa}

La experiencia de usuario constituye el tercer pilar de diseño. Mientras que
herramientas como OWASP ZAP presentan interfaces densas que exigen
especialización previa, y Nikto o Nuclei operan exclusivamente desde la línea
de comandos, VulnSeeker implementa una interfaz gráfica en CustomTkinter con
un paradigma de \textit{un~click, un~escaneo}: el usuario introduce la URL,
selecciona opcionalmente el crawler, y el sistema se encarga del resto. La GUI
incorpora visualizaciones en tiempo real (pie charts de severidad, bar charts
de Top~5), historial con filtros, y exportación directa a PDF. Esta decisión de
diseño responde directamente al perfil de usuario objetivo: desarrolladores y
PYMEs que necesitan resultados rápidos sin formación en ciberseguridad.

\subsection{Memoria de Materiales: Justificación del Stack Tecnológico}

La selección de las tecnologías que conforman VulnSeeker no es arbitraria;
cada componente fue elegido con la deliberación con la que un arquitecto
selecciona los materiales de una obra, evaluando resistencia, compatibilidad
y costo.

\textbf{Python 3.10 o superior} (con entorno de desarrollo bajo versión 3.14) fue seleccionado como lenguaje base por su posición
privilegiada en el ecosistema de ciberseguridad. Así como el acero estructural
combina resistencia y ductilidad, Python combina potencia en scripting de red
con un ecosistema de librerías maduras (\texttt{requests}, \texttt{BeautifulSoup},
\texttt{socket}) que eliminan la necesidad de implementar protocolos desde cero.
Su naturaleza interpretada permite iteraciones rápidas, y su tipado dinámico
facilita la construcción de módulos heterogéneos que comparten una interfaz
común.

\textbf{Llama 3.3 70B} de Meta fue elegido como modelo de inteligencia
artificial por su capacidad de razonamiento lógico y contextualización
técnica. A diferencia de modelos más pequeños, el umbral de 70 mil millones
de parámetros le permite mantener coherencia en informes largos, correlacionar
múltiples hallazgos simultáneamente y generar planes de acción ejecutables. Es
el \textit{concreto armado} del sistema de reportes: soporta la carga
analítica sin ceder ante la complejidad.

\textbf{Groq} fue seleccionado como proveedor de inferencia por su
arquitectura LPU (Language Processing Unit), diseñada específicamente para
inferencia de LLMs. Su latencia inferior a 500ms por consulta convierte al
módulo de IA en una experiencia interactiva, no en un proceso batch. En la
analogía constructiva, Groq es el \textit{sistema de instalaciones} que
permite que la energía (los datos) fluya con mínima resistencia desde el
escaneo hasta el informe ejecutivo.

\textbf{SQLite} como motor de persistencia ofrece lo que en arquitectura se
denominaría \textit{cimentación superficial}: suficiente para la carga
prevista, sin la sobreingeniería de un sistema de base de datos servidor. Su
naturaleza embebida elimina dependencias externas y permite que VulnSeeker
opere como una unidad autónoma desplegable en cualquier equipo con Python
instalado.

\textbf{CustomTkinter} fue seleccionado como framework de interfaz gráfica
siguiendo un criterio de \textit{diseño sostenible y accesibilidad
ergonómica}. En el contexto del software, la sostenibilidad se manifiesta en
la ausencia de dependencias pesadas: CustomTkinter opera sobre Tkinter
---incluido en la biblioteca estándar de Python---, lo cual garantiza que la
GUI funcione sin instalaciones adicionales del sistema operativo, a diferencia
de frameworks como Qt o Electron que arrastran ecosistemas de dependencias de
cientos de megabytes. La accesibilidad ergonómica, por su parte, se traduce en
los widgets nativos con temática oscura que CustomTkinter provee: sliders,
switches y botones con bordes redondeados que no requieren hojas de estilo
adicionales. Para el perfil de usuario objetivo ---desarrolladores y PYMEs sin
formación en UX---, la diferencia entre una interfaz CLI críptica y una GUI
con un botón \textit{``Initialize Scan''} puede significar la diferencia entre
adoptar o abandonar la herramienta.

En conjunto, la selección de materiales responde a un principio unificador:
\textit{máximo rendimiento con mínima dependencia}. Cada componente del stack
fue elegido no por ser el más sofisticado de su categoría, sino por ofrecer la
mejor relación entre capacidad funcional, costo operativo y facilidad de
despliegue ---el equivalente tecnológico de una construcción bioclimática que
aprovecha los recursos disponibles sin exceder su presupuesto energético.

\subsection{Fases Metodológicas}

\subsubsection{Fase 1: Diagnóstico y Análisis (Objetivo Específico 1)}

Se realizó una revisión documental exhaustiva y un benchmarking técnico detallado
de las herramientas DAST y de escaneo de vulnerabilidades de código abierto
existentes en el mercado. En particular, se analizaron de forma comparativa las
capacidades y limitaciones de OWASP~ZAP (versión 2.14), Nikto (versión 2.5),
Wapiti (versión 3.1) y SQLMap (versión 1.8), evaluando criterios como cobertura
funcional frente al OWASP Top~10, extensibilidad arquitectónica, facilidad de
instalación, consumo de recursos del sistema y simplicidad de la interfaz de
usuario.

Esta etapa de diagnóstico permitió constatar tres hallazgos clave: (1) las
herramientas existentes presentan una alta sobrecarga cognitiva en sus interfaces
gráficas o carecen de ellas por completo; (2) la adición de nuevas reglas de
detección exige la comprensión de lógicas complejas en lenguajes compilados o
plantillas declarativas rígidas; y (3) ninguna herramienta integra análisis de
riesgo mediante inteligencia artificial ni consulta automática de vulnerabilidades
conocidas en vivo. A partir de este diagnóstico, se definieron los requisitos
funcionales y no funcionales de VulnSeeker, priorizando la simplicidad ergonómica,
la modularidad pura en Python y el enriquecimiento de reportes mediante LLMs.


\subsubsection{Fase 2: Diseño de la Solución (Objetivo Específico 2)}

Se diseñó una arquitectura modular extensible. El sistema se estructuró en
componentes independientes:

\begin{itemize}
    \item \textbf{Motor de Orquestación (Engine):} Coordina la ejecución
    secuencial y paralela de los módulos de detección.
    \item \textbf{Motor de Mapeo (Crawler):} Descubre y recorre las páginas
    del sitio objetivo, soportando autenticación mediante cookies de sesión.
    \item \textbf{Sistema de Fingerprinting:} Identifica las tecnologías
    subyacentes (servidor, lenguaje, CMS/framework) del objetivo.
    \item \textbf{Gestor de Módulos:} Interfaz base (\texttt{ScannerModule})
    que permite la integración de nuevos detectores sin modificar el núcleo.
    \item \textbf{Motor de Reportes:} Exporta hallazgos en formatos PDF y JSON
    con análisis asistido por inteligencia artificial.
\end{itemize}

El Motor de Orquestación merece una atención particular, pues constituye el
núcleo operativo del framework. Conceptualmente, opera como un
\textit{Centro de Transferencia de Datos Intermodal}: un nodo central que
recibe información heterogénea desde múltiples fuentes de reconocimiento
---URLs descubiertas por el Crawler, tecnologías identificadas por el
Fingerprinter, subdominios revelados por el escáner OSINT--- y la procesa,
normaliza y distribuye eficientemente hacia los 25~módulos de detección.
Análogamente a una terminal logística donde contenedores de distinto origen
convergen para ser clasificados y redirigidos a sus destinos específicos, el
Engine recibe objetos \texttt{Target} con metadatos enriquecidos y los
despacha en paralelo a cada módulo mediante un pool de hilos configurable.
Esta arquitectura intermodal garantiza que ningún módulo necesite realizar
su propio reconocimiento: toda la inteligencia recopilada en la fase previa
fluye automáticamente hacia la fase de detección, eliminando redundancias y
optimizando el tiempo total de escaneo.

\subsubsection{Fase 3: Construcción e Implementación (Objetivo Específico 3)}

Esta fase comprendió la traducción de los diseños lógicos a código de
producción. Se utilizó Python como lenguaje de programación base, aprovechando su
madurez y portabilidad. Los módulos de inyección SQLi y XSS constituyeron el
alcance mínimo inicial del objetivo específico, el cual fue posteriormente ampliado
durante el desarrollo iterativo hasta alcanzar los 25~módulos de la versión final.
La implementación se dividió en dos bloques paralelos de
desarrollo. En el primer bloque se construyó el núcleo operativo del sistema,
programando el motor de orquestación multihilo (\texttt{Engine}), el crawler con
soporte para persistencia de cookies de sesión (\texttt{Crawler}) y el detector de
tecnologías basado en firmas HTTP (\texttt{Fingerprinter}). La persistencia se
resolvió mediante SQLite para almacenar de forma eficiente las URLs descubiertas
y los hallazgos identificados sin requerir dependencias de servidor.

En el segundo bloque se codificaron de forma individual los 25~módulos de
detección, heredando cada uno de la clase abstracta \texttt{ScannerModule} y
siguiendo el patrón Strategy para garantizar que cada detector encapsulara su
propia lógica de inyección y análisis. Adicionalmente, se programaron la interfaz
gráfica de usuario utilizando la librería CustomTkinter ---diseñando widgets
oscuros autoadaptables--- y el conector de inteligencia artificial mediante el SDK de
Groq, estableciendo la estructura del prompt del sistema para orientar la
evaluación del modelo Llama 3.3 70B hacia la accionabilidad técnica.

Los 25~módulos de detección implementados cubren las siguientes categorías del
OWASP Top~10 (2021):


\begin{table}[H]
    \centering
    \caption{Módulos implementados por categoría OWASP}
    \label{tab:modulos}
    \begin{tabular}{p{3cm} p{5.5cm} p{4cm}}
        \toprule
        \textbf{Categoría OWASP} & \textbf{Módulos} & \textbf{Severidad} \\
        \midrule
        A01 — Broken Access Control &
            Open Redirect, CORS Scanner &
            HIGH / MEDIUM \\
        A02 — Crypto Failures &
            TLS Checker, Sensitive Data &
            HIGH / MEDIUM \\
        A03 — Injection &
            SQLi, XSS, Command Injection, LFI, RFI &
            CRITICAL / HIGH \\
        A04 — Insecure Design &
            File Upload Detector &
            HIGH \\
        A05 — Security Misconfig &
            Header Analyzer, Dir Listing, HTTP Methods, Path Fuzzer &
            MEDIUM / HIGH \\
        A06 — Vulnerable Components &
            Tech Fingerprinter, CVE Lookup (NVD API) &
            CRITICAL / HIGH \\
        A07 — Auth Failures &
            CSRF Auditor, Brute Force, Weak Session &
            MEDIUM / HIGH \\
        A10 — SSRF &
            SSRF Scanner &
            HIGH \\
        \bottomrule
    \end{tabular}
\end{table}

\subsubsection{Fase 4: Validación Experimental (Objetivo Específico 4)}

La validación experimental de la solución se estructuró a tres niveles
complementarios para garantizar tanto la calidad del software como su efectividad
diagnóstica. El primer nivel consistió en la construcción de una suite de 150~pruebas
unitarias desarrolladas con el framework \texttt{pytest}. Estas pruebas emplean la
técnica de simulación (\textit{mocking}) a través de \texttt{unittest.mock} para
interceptar las solicitudes de red de cada módulo y retornar respuestas HTTP
prediseñadas, permitiendo verificar el comportamiento de los detectores ante
múltiples escenarios sin depender de conectividad a internet.

El segundo nivel de validación consistió en la automatización del control de
calidad mediante la configuración de un pipeline de Integración Continua (CI/CD)
en GitHub~Actions. Este flujo de trabajo ejecuta automáticamente la suite de tests y
genera reportes de cobertura ante cada cambio en la rama principal de desarrollo,
evitando la introducción de regresiones.

El tercer nivel de validación consistió en la ejecución de auditorías dinámicas
reales en entornos de prueba controlados de la industria: Damn Vulnerable Web
Application (DVWA), desplegada localmente mediante contenedores Docker, y la
plataforma web pública de pruebas Altoro Mutual (demo.testfire.net). Estos escaneos
permitieron medir el comportamiento del crawler ante diferentes estructuras de
enlaces, evaluar la tasa de falsos positivos en entornos heterogéneos y validar la
capacidad de la inteligencia artificial para consolidar y priorizar hallazgos
reales en el reporte final.

% ===========================================================================
% CAPÍTULO IV — RESULTADOS
% ===========================================================================
\newpage
\section{CAPÍTULO IV: RESULTADOS}

\subsection{Resultados de la Validación Unitaria}

La suite de pruebas unitarias, ejecutada mediante el framework \texttt{pytest},
verificó el comportamiento de cada módulo de detección utilizando respuestas
HTTP simuladas (mocking) para garantizar ejecución determinista e independiente
de la red. Los resultados se resumen en la Tabla~\ref{tab:tests}.

\begin{table}[H]
    \centering
    \caption{Resultados de la suite de pruebas unitarias}
    \label{tab:tests}
    \small
    \begin{tabular}{p{4.5cm} c c c}
        \toprule
        \textbf{Módulo} & \textbf{Tests} & \textbf{Passed} & \textbf{Status} \\
        \midrule
        SQL Injection       & 16 & 16 & \checkmark \\
        Cross-Site Scripting & 3  & 3  & \checkmark \\
        Command Injection   & 4  & 4  & \checkmark \\
        Local File Inclusion & 4  & 4  & \checkmark \\
        Remote File Inclusion & 4  & 4  & \checkmark \\
        Open Redirect       & 4  & 4  & \checkmark \\
        CORS Scanner        & 6  & 6  & \checkmark \\
        CSRF Auditor        & 5  & 5  & \checkmark \\
        CVE Lookup          & 9  & 9  & \checkmark \\
        TLS Checker         & 5  & 5  & \checkmark \\
        Header Analyzer     & 3  & 3  & \checkmark \\
        Dir Listing         & 4  & 4  & \checkmark \\
        HTTP Methods        & 5  & 5  & \checkmark \\
        Path Fuzzer         & 5  & 5  & \checkmark \\
        Brute Force         & 6  & 6  & \checkmark \\
        Weak Session        & 6  & 6  & \checkmark \\
        Sensitive Data      & 6  & 6  & \checkmark \\
        File Upload         & 6  & 6  & \checkmark \\
        SSRF Scanner        & 6  & 6  & \checkmark \\
        WAF Detector        & 4  & 4  & \checkmark \\
        CMS Auditor         & 4  & 4  & \checkmark \\
        Cookie Scanner      & 4  & 4  & \checkmark \\
        Email Harvester     & 4  & 4  & \checkmark \\
        Subdomain Takeover  & 4  & 4  & \checkmark \\
        Port Scanner        & 2  & 2  & \checkmark \\
        Auto-login DVWA (core) & 4 & 4 & \checkmark \\
        Models (core)       & 5  & 5  & \checkmark \\
        Engine (core)       & 5  & 5  & \checkmark \\
        Crawler (core)      & 7  & 7  & \checkmark \\
        Fingerprinter (core) & 8 & 8  & \checkmark \\
        API Discovery (core) & 6 & 6  & \checkmark \\
        \midrule
        \textbf{TOTAL}      & \textbf{164} & \textbf{164} & \textbf{100\%} \\
        \bottomrule
    \end{tabular}
\end{table}

El 100\% de los tests pasaron exitosamente. La integración continua mediante
GitHub Actions garantiza que cualquier regresión sea detectada automáticamente
en cada commit.

\subsection{Resultados de los Escaneos Controlados}

Se ejecutó un escaneo completo con los 25~módulos activos contra DVWA
(Damn~Vulnerable Web~Application) desplegada localmente en un contenedor
Docker con nivel de seguridad configurado en \textit{``low''} para maximizar
la superficie de ataque detectable. El crawler autenticado ---tras realizar
el inicio de sesión de forma automática mediante la extracción del token CSRF
de DVWA--- identificó 61~puntos de entrada estructuralmente únicos (endpoints
y formularios), sobre los cuales cada módulo ejecutó su batería de pruebas.

\begin{table}[H]
    \centering
    \caption{Resumen del escaneo de validación contra DVWA}
    \label{tab:scan-dvwa}
    \begin{tabular}{l r}
        \toprule
        \textbf{Métrica} & \textbf{Valor} \\
        \midrule
        Total de hallazgos        & 395 \\
        Críticos (CRITICAL)       & 3 (0,8\%) \\
        Altos (HIGH)              & 61 (15,4\%) \\
        Medios (MEDIUM)           & 75 (19,0\%) \\
        Bajos (LOW)               & 255 (64,6\%) \\
        Informativos (INFO)       & 1 (0,3\%) \\
        Puntos de entrada únicos  & 61 \\
        Subdominios descubiertos  & 0 (escaneo local) \\
        Risk Score (IA)           & 100 / 100 --- CRÍTICO \\
        \bottomrule
    \end{tabular}
\end{table}

Buena parte de los hallazgos corresponde a deficiencias de configuración
evaluadas de forma independiente en cada endpoint: por ejemplo, cada punto de
entrada servido sobre HTTP sin el header \texttt{X-Frame-Options} genera un
hallazgo individual, lo cual es coherente con el enfoque DAST de auditoría
exhaustiva por endpoint. El motor aplica una etapa de deduplicación
post-escaneo que elimina hallazgos exactamente idénticos (mismo nombre de
vulnerabilidad, URL objetivo y payload), conservando intencionalmente los del
mismo tipo en URLs distintas, ya que cada endpoint representa un punto de
entrada independiente que requiere remediación individual.

\subsubsection{Depuración de Falsos Positivos}

Un resultado metodológicamente relevante del proceso de validación fue la
reducción sistemática de la tasa de falsos positivos a lo largo de la
evolución del framework. Las versiones iniciales de los módulos de inyección
activa carecían de un mecanismo de comparación contra una respuesta de
referencia (\textit{baseline}): un mismo indicador ---por ejemplo, la cadena
esperada de la salida de un comando del sistema--- se reportaba en numerosos
endpoints donde ya estaba presente de forma legítima, inflando artificialmente
el número de hallazgos. La introducción de la validación contra baseline
---verificar que el indicador no estuviera presente en la respuesta original
\textit{antes} del ataque--- depuró estos falsos positivos. La
Tabla~\ref{tab:fp} ilustra esta evolución para el módulo de Command Injection,
cuyo conteo de hallazgos en DVWA ---un entorno con un único punto real de
inyección de comandos--- pasó de 449 (mayoritariamente falsos positivos) a un
único verdadero positivo.

\begin{table}[H]
    \centering
    \caption{Reducción de falsos positivos de Command Injection (RCE) en DVWA}
    \label{tab:fp}
    \begin{tabular}{l c c}
        \toprule
        \textbf{Versión del framework} & \textbf{Hallazgos RCE} & \textbf{Naturaleza} \\
        \midrule
        Inicial (sin baseline)     & 449 & Mayoría falsos positivos \\
        Validación parcial         & 18  & Falsos positivos residuales \\
        Final (baseline completo)  & 1   & Verdadero positivo (\texttt{/exec/}) \\
        \bottomrule
    \end{tabular}
\end{table}

Esta depuración constituye una contribución concreta del trabajo: prioriza la
\textit{precisión} (ausencia de falsos positivos) sobre el volumen bruto de
hallazgos, en coherencia con el criterio de minimización de la tasa de falsos
positivos descrito en el marco teórico y con la realidad operativa de los
equipos de seguridad, que padecen fatiga de alertas ante reportes ruidosos.

\subsection{Top~5 Vulnerabilidades Recurrentes}

El análisis de frecuencia reveló que las vulnerabilidades más recurrentes
corresponden a deficiencias de configuración de seguridad (OWASP~A05),
detectadas de forma uniforme en los endpoints servidos por la aplicación
(55~ocurrencias cada una):

\begin{table}[H]
    \centering
    \caption{Top~5 vulnerabilidades recurrentes en DVWA}
    \label{tab:top5}
    \begin{tabular}{c l c c}
        \toprule
        \textbf{\#} & \textbf{Vulnerabilidad} & \textbf{Ocurrencias} &
        \textbf{Severidad} \\
        \midrule
        1 & No HTTPS                         & 55 & HIGH \\
        2 & Missing HSTS                     & 55 & HIGH \\
        3 & Missing X-Frame-Options          & 55 & MEDIUM \\
        4 & Missing Content-Security-Policy   & 55 & MEDIUM \\
        5 & Missing X-Content-Type-Options    & 55 & MEDIUM \\
        \bottomrule
    \end{tabular}
\end{table}

Estos resultados son consistentes con el comportamiento esperado: DVWA es una
aplicación intencionalmente vulnerable desplegada sobre HTTP sin configuración
de headers de seguridad. La detección uniforme en los endpoints servidos
confirma que el crawler recorrió la superficie de la aplicación y que el
módulo Header~Analyzer evaluó cada endpoint de forma independiente.

\subsection{Efectividad por Módulo}

Para evaluar la precisión de detección, los hallazgos reportados por cada
módulo fueron sometidos a verificación manual cruzada contra las
vulnerabilidades documentadas de DVWA, con el fin de contabilizar los
verdaderos positivos (VP) y los falsos positivos (FP) observados en el
escenario experimental. La Tabla~\ref{tab:efectividad} resume estos valores.

\begin{table}[H]
    \centering
    \caption{Precisión observada por módulo contra DVWA}
    \label{tab:efectividad}
    \small
    \begin{tabular}{p{4cm} p{4.2cm} c c}
        \toprule
        \textbf{Módulo} & \textbf{VP (ubicación en DVWA)} & \textbf{FP} &
        \textbf{Precisión} \\
        \midrule
        \multicolumn{4}{l}{\textit{Módulos de inyección activa}} \\
        SQLi (error-based)  & 2 (\texttt{/sqli/}, \texttt{/brute/}) & 0 & 100\% \\
        SQLi (blind time-based) & 1 (\texttt{/sqli\_blind/})    & 0 & 100\% \\
        Reflected XSS      & 1 (\texttt{/xss\_r/})             & 0 & 100\% \\
        Command Injection  & 1 (\texttt{/exec/})               & 0 & 100\% \\
        LFI                & 1 (\texttt{/fi/})                 & 0 & 100\% \\
        RFI                & 1 (\texttt{/fi/})                 & 0 & 100\% \\
        File Upload        & 2 (\texttt{/upload/})             & 0 & 100\% \\
        Open Redirect      & 1 (\texttt{/open\_redirect/})     & 0 & 100\% \\
        \midrule
        \multicolumn{4}{l}{\textit{Módulos de configuración (1 hallazgo por endpoint)}} \\
        Header Analyzer    & 220 (4 headers $\times$ 55)       & 0 & 100\% \\
        TLS / No HTTPS     & 55                                & 0 & 100\% \\
        Server Disclosure  & 53                                & 0 & 100\% \\
        CSRF Auditor       & 9                                 & 0 & 100\% \\
        \bottomrule
    \end{tabular}
\end{table}

Los módulos de inyección activa detectaron correctamente cada clase de
vulnerabilidad en su ubicación esperada dentro de DVWA ---SQL Injection
basada en errores en los formularios de \texttt{/sqli/} y \texttt{/brute/},
SQL Injection ciega basada en tiempo en \texttt{/sqli\_blind/}, XSS reflejado
en \texttt{/xss\_r/}, ejecución de comandos en \texttt{/exec/}, inclusión de
archivos local y remota en \texttt{/fi/}, carga de archivos en \texttt{/upload/}
y redirección abierta en \texttt{/open\_redirect/}--- sin generar falsos
positivos observados, gracias al uso de payloads específicos, al análisis de
patrones de error y de latencia, y a la validación contra baseline descrita en
la sección anterior. Los módulos de configuración reportaron un hallazgo por endpoint
afectado, lo cual constituye una detección verdaderamente positiva dado que
DVWA opera sobre HTTP sin headers de seguridad configurados.

Es importante precisar el alcance de estas métricas. La precisión observada
($VP / (VP + FP)$) cuantifica la ausencia de falsos positivos en los hallazgos
verificados manualmente, y constituye el indicador priorizado por el diseño de
VulnSeeker. En cambio, el cálculo formal de la
sensibilidad o \textit{recall} ($VP / (VP + FN)$) exige un conjunto de prueba
con \textit{ground truth} exhaustivamente etiquetado ---es decir, el inventario
completo y verificado de todas las vulnerabilidades existentes en el objetivo---,
condición que DVWA no garantiza por sí sola. Por esta razón, el presente trabajo
reporta la precisión observada y reserva la medición formal del recall y del
Índice de Youden para una validación contra el OWASP Benchmark Project, planteada
como trabajo futuro (véase la sección de Recomendaciones). Esta decisión
metodológica evita sobreestimar la cobertura real del framework, en coherencia
con las limitaciones reconocidas más adelante.

\subsection{Resultados del Escaneo Externo contra Altoro Mutual (Scan~\#93)}

Con el propósito de evaluar la efectividad del framework en infraestructuras desplegadas en la internet pública, se ejecutó una auditoría dinámica
(Scan~\#93) contra la plataforma de pruebas Altoro Mutual (\texttt{demo.testfire.net}).
A diferencia de DVWA, que opera bajo una red local Dockerizada y sin latencia, este escaneo
externo permitió verificar el comportamiento del crawler ante respuestas dinámicas,
múltiples parámetros y latencia de red variable.

El escaneo activo identificó un total de 268~hallazgos válidos distribuidos en diferentes
niveles de severidad, como se resume en la Tabla~\ref{tab:scan-altoro}.

\begin{table}[H]
    \centering
    \caption{Distribución de hallazgos por severidad en Altoro Mutual (Scan~\#93)}
    \label{tab:scan-altoro}
    \begin{tabular}{l c c}
        \toprule
        \textbf{Severidad} & \textbf{Ocurrencias} & \textbf{Porcentaje} \\
        \midrule
        Alto (HIGH)          & 18  & 6,7\% \\
        Medio (MEDIUM)       & 95  & 35,4\% \\
        Bajo (LOW)           & 155 & 57,8\% \\
        Informativo (INFO)   & 0   & 0,0\% \\
        \midrule
        \textbf{TOTAL}       & \textbf{268} & \textbf{100,0\%} \\
        \bottomrule
    \end{tabular}
\end{table}

El análisis detallado de los hallazgos reveló debilidades de configuración crítica
y fallos de inyección activa:

\begin{itemize}
    \item \textbf{Métodos HTTP Peligrosos (HIGH):} El módulo HTTP Method Scanner
    detectó que el servidor tiene habilitados los métodos \texttt{PUT} (9 ocurrencias)
    y \texttt{DELETE} (9 ocurrencias) en los directorios de documentación de la API. Esta
    configuración permite a un atacante no autenticado subir archivos arbitrarios o eliminar
    recursos en el servidor, representando un riesgo crítico de intrusión.
    
    \item \textbf{Reflected Cross-Site Scripting (XSS) (MEDIUM):} El módulo XSS detectó
    que el parámetro \texttt{content} en la página \texttt{index.jsp} refleja las entradas
    del usuario sin saneamiento alguno. La inyección del payload \texttt{<VulnSeekerXSS>} fue
    recibida de vuelta íntegramente en la respuesta HTTP, confirmando la vulnerabilidad.

    \item \textbf{Ausencia de Atributo SameSite en Cookies (MEDIUM):} El módulo Cookie
    Scanner reportó 25~ocurrencias de cookies de sesión (\texttt{JSESSIONID}) que carecen del
    atributo \texttt{SameSite}. Al no contar con esta protección, el navegador envía la cookie en
    solicitudes cruzadas, lo cual expone a los usuarios de la aplicación a ataques de
    Cross-Site Request Forgery (CSRF).

    \item \textbf{Exposición de Puertos de Administración (MEDIUM):} El módulo Port Scanner
    reportó que el puerto alternativo HTTP \texttt{8080} se encuentra expuesto públicamente,
    lo que permite deducir que la consola de administración del servidor de aplicaciones
    (Apache Tomcat / Coyote) es accesible desde el exterior.

    \item \textbf{Deficiencias de Configuración de Cabeceras (LOW):} Se identificaron
    31~ocurrencias de omisión de cabeceras de seguridad fundamentales, incluyendo
    \texttt{X-Frame-Options} (Clickjacking), \texttt{Content-Security-Policy} (ejecución de XSS),
    \texttt{Strict-Transport-Security} (downgrade HTTPS) y \texttt{X-Content-Type-Options}
    (confusión de tipos MIME). Estas deficiencias se detectaron uniformemente en las
    31~páginas del portal indexadas por el crawler.
\end{itemize}

Este escaneo externo valida la robustez de VulnSeeker ante variables operativas de red real: la herramienta
completó la auditoría de forma remota bajo condiciones normales de tráfico y latencia
sin generar denegación de servicio (DoS) ni bloqueos de red, demostrando la precisión de sus
motores heurísticos en entornos no locales.

\subsection{Validación contra una SPA: OWASP Juice Shop y la Extensión API-Aware}

Las aplicaciones de página única (\textit{Single-Page Applications}, SPA)
construidas con frameworks como Angular, React o Vue representan un reto
estructural para los escáneres DAST basados en \textit{crawling} estático. En
estas arquitecturas, la navegación, los formularios y las llamadas al
\textit{backend} se generan dinámicamente mediante JavaScript en el cliente y se
comunican con una API REST; el HTML inicial apenas contiene un contenedor raíz
vacío. Para verificar el comportamiento de VulnSeeker en este escenario se
desplegó localmente OWASP~Juice~Shop (Node.js/Express + Angular), considerada la
aplicación insegura moderna de referencia.

\subsubsection{El problema de visibilidad y su solución sin navegador}

El crawler estructural, al no ejecutar JavaScript, mapeó 0~endpoints sobre la
SPA: el motor no disponía de superficie sobre la cual atacar. Antes que recurrir
a un motor de renderizado pesado (Selenium/Playwright) ---con su coste de
dependencias, fragilidad y consumo de recursos--- se incorporó un módulo de
\textbf{descubrimiento de endpoints API sin navegador} que combina tres capas:
(i)~el parseo de una especificación OpenAPI/Swagger cuando la aplicación la
publica; (ii)~la extracción por expresiones regulares de las rutas \texttt{/api/}
y \texttt{/rest/} embebidas en los \textit{bundles} JavaScript de la aplicación;
y (iii)~el sondeo validado de \textit{endpoints} REST canónicos (búsqueda y
autenticación). El efecto sobre la visibilidad fue determinante
(Tabla~\ref{tab:spa-visibilidad}).

\begin{table}[H]
    \centering
    \caption{Visibilidad de la superficie de ataque en OWASP Juice Shop}
    \label{tab:spa-visibilidad}
    \begin{tabular}{l c}
        \toprule
        \textbf{Mecanismo de descubrimiento} & \textbf{Endpoints mapeados} \\
        \midrule
        Crawler estructural (estático)      & 0 \\
        Descubridor API-aware (spec + JS)   & 29 \\
        \bottomrule
    \end{tabular}
\end{table}

\subsubsection{Robustez ante respuestas \textit{catch-all} (soft-404)}

El despliegue de Juice Shop reveló un patrón que los entornos previos (DVWA,
Altoro) no exhibían: como toda SPA, el servidor responde \texttt{HTTP~200} con su
\texttt{index.html} ante \textit{cualquier} ruta inexistente (comportamiento
\textit{catch-all} o \textit{soft-404}), para que el enrutado se resuelva en el
cliente. Los módulos que infieren la existencia de un recurso a partir de un
código 200 ---como el descubridor de archivos sensibles o la huella tecnológica
por rutas--- generaban en consecuencia falsos positivos masivos: reportaban como
expuestos 24~archivos inexistentes (\texttt{.env}, \texttt{.git/config},
\texttt{wp-config.php.bak}, copias de base de datos, etc.) e identificaban
erróneamente el \textit{stack} como ``Joomla/WordPress''. Se introdujo un
mecanismo de \textit{baseline} que aprende la firma de la página
\textit{catch-all} ---sondeando rutas aleatorias garantizadamente inexistentes---
y descarta toda respuesta 200 que coincida con ella (Tabla~\ref{tab:soft404}). En
servidores con 404 real (DVWA) el mecanismo no se activa y el comportamiento del
escáner no varía.

\begin{table}[H]
    \centering
    \caption{Depuración de falsos positivos por respuesta catch-all en Juice Shop}
    \label{tab:soft404}
    \begin{tabular}{l c c}
        \toprule
        \textbf{Indicador} & \textbf{Sin baseline} & \textbf{Con baseline} \\
        \midrule
        Archivos sensibles ``expuestos'' (FP) & 24 & 0 \\
        Huella de CMS falsa (Joomla/WordPress) & Sí & No \\
        \bottomrule
    \end{tabular}
\end{table}

\subsubsection{Hallazgos y verdaderos positivos de inyección}

Con la superficie de API restaurada, el escaneo completo produjo 117~hallazgos.
Conviene interpretarlos con rigor: la mayoría son deficiencias de configuración
(ausencia de HTTPS, CORS con origen comodín, ausencia de CSP y HSTS) evaluadas de
forma independiente en cada uno de los 29~endpoints, por lo que no constituyen
117~problemas distintos sino aproximadamente siete tipos replicados por endpoint.
Lo relevante a efectos de validación es que el motor de inyección, extendido para
operar sobre cuerpos JSON y dotado de firmas de error de SQLite y de ORM
(Sequelize), confirmó dos~vulnerabilidades de inyección SQL reales
(Tabla~\ref{tab:juiceshop-sqli}): una en un parámetro de consulta (GET) y otra en
el cuerpo JSON del \textit{endpoint} de autenticación ---donde el payload
\texttt{' OR 1=1--} produce además un \textit{bypass} de credenciales con emisión
de un token JWT de administrador.

\begin{table}[H]
    \centering
    \caption{Verdaderos positivos de SQL Injection detectados en Juice Shop}
    \label{tab:juiceshop-sqli}
    \small
    \begin{tabular}{p{3.8cm} p{2.4cm} p{5.0cm}}
        \toprule
        \textbf{Endpoint} & \textbf{Vector} & \textbf{Evidencia} \\
        \midrule
        \texttt{/rest/products/search} & Query GET (\texttt{q}) & Error \texttt{SQLITE\_ERROR} ante payload de ruptura \\
        \texttt{/rest/user/login} & Cuerpo JSON (\texttt{email}) & \textit{Stack} de \texttt{sequelize/dialects/sqlite}; bypass con \texttt{' OR 1=1--} \\
        \bottomrule
    \end{tabular}
\end{table}

Este caso aporta una doble lección metodológica. Por un lado, delimita la
\textbf{validez externa} del framework: VulnSeeker cubre con solvencia las
aplicaciones \textit{server-rendered} (DVWA, Altoro), mientras que en las SPA su
eficacia depende de la capa de descubrimiento API-aware aquí incorporada, que
recupera la superficie de ataque sin el coste de un navegador \textit{headless}.
Por otro, confirma que el endurecimiento frente a falsos positivos (\textit{baseline}
\textit{catch-all}) es indispensable para operar sobre aplicaciones modernas. Como
trabajo futuro queda la integración de un motor de renderizado \textit{headless}
para SPAs que no publiquen especificación ni expongan sus rutas de forma legible
en el código JavaScript.

\subsection{Evaluación del Análisis de Inteligencia Artificial}


El módulo de análisis IA, basado en Llama 3.3 70B vía Groq API, fue evaluado
cualitativamente bajo cinco criterios:

\begin{table}[H]
    \centering
    \caption{Evaluación cualitativa del análisis de IA}
    \label{tab:ia-eval}
    \begin{tabular}{p{4cm} p{2cm} p{6.5cm}}
        \toprule
        \textbf{Criterio} & \textbf{Resultado} & \textbf{Observación} \\
        \midrule
        Relevancia &
            Alta &
            El modelo identifica correctamente las vulnerabilidades más
            críticas del escaneo. \\
        Coherencia &
            Alta &
            Las recomendaciones son técnicamente consistentes con los
            hallazgos reportados. \\
        Accionabilidad &
            Alta &
            El plan de acción inmediato (48h) incluye pasos específicos
            y priorizados. \\
        Lenguaje ejecutivo &
            Alta &
            El informe traduce terminología técnica a riesgo de negocio
            comprensible. \\
        Alucinaciones &
            Baja &
            No se detectaron fabricaciones de vulnerabilidades
            inexistentes en las pruebas realizadas. \\
        \bottomrule
    \end{tabular}
\end{table}

\subsection{Análisis e Interpretación de Resultados}

Los resultados obtenidos permiten una interpretación a tres niveles:
la validación de la hipótesis arquitectónica, la confrontación con los
antecedentes de investigación y las implicaciones prácticas para el
ecosistema objetivo.

\subsubsection{Validación de la Hipótesis Arquitectónica}

La decisión de implementar una arquitectura modular basada en el patrón
Strategy ---donde cada módulo de detección opera como una estrategia
intercambiable según Gamma et~al. (\citeyear*{gamma1995})--- se valida
por los resultados de la suite de pruebas unitarias: los 150~tests
pasaron con un 100\% de aprobación, demostrando que cada módulo opera
de forma independiente sin generar regresiones en los demás. Esta
independencia confirma experimentalmente lo que Bass et~al.
(\citeyear*{bass2012}) denominan el atributo de \textit{testabilidad}:
la capacidad de un sistema para facilitar la verificación aislada de
cada uno de sus componentes.

La adición de nuevos módulos durante el desarrollo ---desde los 2
módulos iniciales (SQLi, XSS) hasta los 25~finales--- se realizó sin
modificar una sola línea del motor de orquestación, validando el
principio Open/Closed de Martin (\citeyear*{martin2003}) y confirmando
la viabilidad del prototipado evolutivo de Sommerville
(\citeyear*{sommerville2011}) aplicado a la construcción de herramientas
de seguridad.

\subsubsection{Confrontación con los Antecedentes}

Los resultados de VulnSeeker son consistentes con los hallazgos de la
literatura revisada:

\begin{itemize}
    \item \textbf{Abdulghaffar et al. (\citeyear*{abdulghaffar2023})}
    demostraron que la fusión de múltiples motores de detección supera
    la tasa de detección individual. VulnSeeker confirma esta tesis:
    los 25~módulos operando en paralelo detectaron vulnerabilidades
    en 8 de las 10~categorías del OWASP Top~10, una cobertura que
    ningún módulo individual podría alcanzar.

    \item \textbf{Alhogail y Alkahtani (\citeyear*{alhogail2024})}
    validaron que la extensibilidad basada en extensiones mejora la
    escalabilidad sin alterar el núcleo. La arquitectura de VulnSeeker
    reproduce este hallazgo: la clase \texttt{ScannerModule} permite
    la incorporación de nuevos detectores mediante herencia simple,
    sin modificar el engine ni el registro de módulos.

    \item \textbf{Yarphel y Rani (\citeyear*{yarphel2025})} concluyeron
    que ninguna herramienta cubre todas las vulnerabilidades por sí sola.
    VulnSeeker no contradice esta conclusión ---las categorías A08 y A09
    escapan al alcance de cualquier herramienta DAST---, pero demuestra
    que un framework modular puede maximizar la cobertura dentro de los
    límites inherentes a la metodología de caja negra.

    \item \textbf{Repudi y Harinadh (\citeyear*{repudi2024})} reportaron
    una reducción del 40\% en tiempos de escaneo mediante procesamiento
    multihilo. VulnSeeker implementa un pool de hilos configurable que
    ejecuta los 25~módulos en paralelo, confirmando la viabilidad del
    diseño concurrente para herramientas DAST.
\end{itemize}

\subsubsection{Implicaciones Prácticas}

Desde la perspectiva del ecosistema objetivo ---desarrolladores
independientes y PYMEs---, los resultados tienen tres implicaciones
principales:

\begin{enumerate}
    \item \textbf{Costo incremental de licenciamiento nulo:} Como se detalla en la
    Tabla~\ref{tab:costos}, VulnSeeker opera íntegramente con
    tecnologías gratuitas, eliminando la barrera económica identificada
    en el planteamiento del problema y confirmada por Altulaihan et~al.
    (\citeyear*{altulaihan2023}) como una de las principales limitantes
    del ecosistema actual.

    \item \textbf{Reducción de la brecha regional:} En el contexto
    del informe OEA/BID (\cite{oasidb2025}), que documenta las brechas
    de madurez en ciberseguridad en América Latina, VulnSeeker contribuye
    directamente a la dimensión de \textit{herramientas accesibles} que
    el informe identifica como área de mejora prioritaria.

    \item \textbf{Mitigación del riesgo de código IA:} Dado que entre
    el 45\% y el 48\% del código generado por IA contiene fallos de
    seguridad, VulnSeeker ofrece una capa de validación post-desarrollo
    que puede detectar las vulnerabilidades más comunes introducidas por
    asistentes de programación automatizados.
\end{enumerate}

% ===========================================================================
% CAPÍTULO V — CONCLUSIONES Y RECOMENDACIONES
% ===========================================================================
\newpage
\section{CAPÍTULO V: CONCLUSIONES Y RECOMENDACIONES}

\subsection{Conclusiones}

A partir de los resultados obtenidos y en correspondencia con los objetivos
específicos planteados, se derivan las siguientes conclusiones:

\begin{enumerate}
    \item \textbf{Objetivo 1 — Análisis de herramientas existentes:} El estudio
    comparativo de OWASP ZAP, Nikto, Nuclei y Burp Suite reveló que ninguna
    herramienta de código abierto integra simultáneamente análisis DAST
    modular, consulta de CVEs en tiempo real e informes asistidos por
    inteligencia artificial. Este hallazgo es consistente con la revisión
    sistemática de Altulaihan et~al. (\citeyear*{altulaihan2023}), quienes
    documentaron que las herramientas actuales suelen carecer de modularidad
    o poseen costos prohibitivos. La carencia identificada validó la
    pertinencia de desarrollar VulnSeeker como solución integrada.

    \item \textbf{Objetivo 2 — Diseño de la arquitectura:} La arquitectura
    modular implementada, basada en los principios SOLID de Martin
    (\citeyear*{martin2003}), el patrón Strategy de Gamma et~al.
    (\citeyear*{gamma1995}) y la Regla de Dependencia de la arquitectura
    limpia (\cite{martin2017}), demostró ser extensible sin modificar el
    núcleo. Los 25~módulos operan de forma independiente y se registran
    dinámicamente en el motor de orquestación, confirmando los hallazgos
    de Alhogail y Alkahtani (\citeyear*{alhogail2024}) sobre la viabilidad
    del enfoque extensible para herramientas de pentesting.

    \item \textbf{Objetivo 3 — Implementación de módulos:} Se implementaron
    exitosamente 25~módulos de detección que cubren 8 de las 10~categorías
    del OWASP Top~10 (2021). La integración con la NVD API~2.0 del NIST y
    el modelo Llama 3.3 70B amplió las capacidades de la herramienta más allá
    de la detección estática, integrando análisis asistido por inteligencia artificial
    generativa dentro de una arquitectura DAST modular.

    \item \textbf{Objetivo 4 — Validación experimental:} Las 150~pruebas
    unitarias alcanzaron un 100\% de aprobación. Los escaneos controlados
    contra DVWA confirmaron la detección correcta de las distintas clases de
    vulnerabilidad (SQLi, XSS, Command Injection, LFI, RFI, File Upload y Open
    Redirect), cada una en su ubicación esperada y sin falsos positivos
    observados en los módulos de inyección. Adicionalmente, la introducción de
    la validación contra baseline depuró los falsos positivos de inyección
    (reduciendo, por ejemplo, los hallazgos de Command Injection de 449 a un
    único verdadero positivo), priorizando la precisión sobre el volumen bruto.
    Estos resultados validan la tesis central de Abdulghaffar et~al.
    (\citeyear*{abdulghaffar2023}): la integración de múltiples motores en una
    arquitectura concurrente amplía la cobertura, abarcando 8~categorías del
    OWASP Top~10.
\end{enumerate}

Se concluye que la integración de múltiples motores de detección en una
arquitectura modular extensible incrementa la cobertura de vulnerabilidades
detectadas sin sacrificar la mantenibilidad del sistema. VulnSeeker demuestra
que es viable construir una herramienta DAST competitiva y accesible
utilizando exclusivamente tecnologías de código abierto, validando tanto la
hipótesis técnica como la viabilidad económica planteada en el Capítulo~I.

\subsection{Impacto en el Paisaje de Seguridad Digital}

Más allá de su contribución técnica, VulnSeeker incide en el paisaje de
seguridad digital de las PYMEs venezolanas. En un contexto donde las
restricciones económicas y las limitaciones de acceso a licencias comerciales
han generado un ecosistema de aplicaciones web estructuralmente vulnerables,
la disponibilidad de una herramienta de auditoría gratuita, modular y operada
localmente representa un vector de \textit{resiliencia digital}: la capacidad
de una organización para anticipar, resistir y recuperarse de incidentes de
ciberseguridad. Al democratizar el acceso a la evaluación dinámica de
vulnerabilidades, VulnSeeker contribuye a reducir la asimetría de seguridad
entre grandes corporaciones ---que disponen de equipos de Red Team y licencias
comerciales--- y los emprendimientos locales que sostienen el tejido económico
nacional. En última instancia, cada PYME que incorpore una herramienta como
VulnSeeker en su ciclo de desarrollo no solo protege sus activos digitales,
sino que fortalece la confianza del usuario final en el ecosistema digital
venezolano.

\subsection{Limitaciones del Estudio}

Es pertinente reconocer las limitaciones inherentes al presente trabajo,
tanto para contextualizar los resultados como para orientar futuras
investigaciones:

\begin{enumerate}
    \item \textbf{Detección heurística:} Los módulos de inyección combinan
    técnicas basadas en errores (\textit{error-based}), análisis de respuestas
    HTTP y ---en el caso de SQL Injection--- detección ciega basada en tiempo
    (\textit{time-based}) mediante la inyección de retardos controlados con
    confirmación proporcional. No obstante, otras variantes avanzadas como el
    XSS almacenado (\textit{stored}) y el DOM XSS no son detectadas en la
    versión actual, por requerir el seguimiento de estado entre peticiones o la
    ejecución de JavaScript del lado del cliente. Esta es una limitación
    compartida con la mayoría de escáneres DAST de código abierto.

    \item \textbf{SSRF in-band:} El módulo SSRF detecta vulnerabilidades
    mediante análisis heurístico de la respuesta HTTP (presencia de
    contenido interno o cambio significativo de tamaño). No utiliza
    técnicas \textit{out-of-band} (OOB) con servidores de callback
    externos, lo cual limita la detección de SSRF ciego.

    \item \textbf{Subdominios como reconocimiento:} El módulo de
    descubrimiento de subdominios opera exclusivamente como OSINT
    informativo. Los subdominios descubiertos no se incorporan
    automáticamente al crawler como URLs de ataque, ya que el escaneo
    activo de subdominios requiere autorización explícita del operador
    para cada dominio.

    \item \textbf{Dependencia de API externa para IA:} El módulo de
    análisis con inteligencia artificial depende de la API de Groq para
    la inferencia con Llama 3.3 70B. En ausencia de conectividad o ante
    cambios en el servicio, los reportes ejecutivos no están disponibles,
    aunque el resto del framework opera con normalidad.

    \item \textbf{Entorno de validación y cobertura de SPAs:} La medición de
    precisión (0~falsos positivos observados) se obtuvo contra DVWA, un entorno
    intencionalmente vulnerable, y se complementó con escaneos sobre la plataforma
    pública Altoro Mutual y sobre la SPA OWASP~Juice~Shop. Esta última evidenció
    que el crawler estático no alcanza la superficie de las aplicaciones que
    construyen su API en JavaScript del lado del cliente; la limitación se mitigó
    incorporando un descubridor de \textit{endpoints} API-aware sin navegador
    (especificación OpenAPI + extracción de rutas en los \textit{bundles} JS), si
    bien las SPAs que no publiquen especificación ni expongan sus rutas de forma
    legible requerirían un motor de renderizado \textit{headless}, lo que queda
    como trabajo futuro. Del mismo modo, la cuantificación del \textit{recall} y
    del Índice de Youden contra un conjunto con \textit{ground truth} etiquetado
    (p.~ej. OWASP~Benchmark) permanece como línea de trabajo pendiente.
\end{enumerate}

\subsection{Recomendaciones}

A partir del trabajo realizado, se formulan las siguientes recomendaciones
para futuras investigaciones y desarrollo:

\begin{enumerate}
    \item Implementar un modo de escaneo asíncrono con \texttt{asyncio} para
    mejorar el rendimiento en auditorías de gran escala con cientos de
    endpoints.

    \item Incorporar un motor de renderizado JavaScript (Playwright o
    Selenium) para evaluar aplicaciones SPA (Single Page Applications)
    que generan contenido dinámico.

    \item Expandir la cobertura del OWASP Top~10 mediante la implementación
    de módulos experimentales para las categorías A08 y A09, utilizando
    heurísticas indirectas.

    \item Desarrollar un sistema de correlación automática de hallazgos
    (Attack Chain Analysis) que permita a la IA inferir cadenas de ataque
    combinando múltiples vulnerabilidades individuales.

    \item Evaluar la efectividad de VulnSeeker contra entornos adicionales
    como OWASP Juice Shop y aplicaciones web reales (con autorización) para
    ampliar la validez externa de los resultados.

    \item Ejecutar una validación formal de la suite completa de detección
    contra el estándar del OWASP Benchmark Project para obtener el Índice de
    Youden exacto del framework, permitiendo la comparación objetiva de
    VulnSeeker con otros motores comerciales.

    \item Diseñar adaptadores de integración continua que faciliten la
    incorporación automática de VulnSeeker en pipelines de desarrollo (tales
    como GitHub Actions, GitLab CI/CD o Jenkins), habilitando la ejecución
    de escaneos dinámicos automáticos en cada despliegue de desarrollo
    (DevSecOps).
\end{enumerate}


\subsection{Memoria Descriptiva de la Solución}

\subsubsection{Arquitectura del Sistema}

VulnSeeker se estructura en cinco capas funcionales que operan de forma
desacoplada:

\begin{enumerate}
    \item \textbf{Capa de Entrada:} Interfaz gráfica (GUI) y línea de
    comandos (CLI) que reciben el objetivo y los parámetros de configuración.

    \item \textbf{Capa de Reconocimiento:} Crawler autenticado,
    Fingerprinting tecnológico y escáner OSINT de subdominios.

    \item \textbf{Capa de Detección:} 25~módulos independientes coordinados
    por el motor de orquestación mediante ejecución multihilo.

    \item \textbf{Capa de Inteligencia:} Consulta a NVD API para CVEs y
    análisis IA para informes ejecutivos.

    \item \textbf{Capa de Salida:} Generación de reportes en PDF, JSON y
    CSV, con persistencia en SQLite.
\end{enumerate}

\subsubsection{Análisis de Costos Operativos}

\begin{table}[H]
    \centering
    \caption{Análisis de costos operativos de VulnSeeker}
    \label{tab:costos}
    \begin{tabular}{p{4cm} p{3cm} p{5cm}}
        \toprule
        \textbf{Recurso} & \textbf{Costo} & \textbf{Observación} \\
        \midrule
        Licencia del software &
            \$0 (MIT) &
            Código abierto, sin restricciones de uso \\
        API de Groq (IA) &
            Gratis (tier free) &
            Límite de 30 RPM en plan gratuito \\
        NVD API del NIST &
            Gratis &
            Acceso público sin clave requerida \\
        Infraestructura &
            \$0 &
            Ejecutable en cualquier PC con Python 3.10+ \\
        CI/CD (GitHub Actions) &
            Gratis &
            2.000 minutos/mes en plan free \\
        \midrule
        \textbf{Costo incremental} &
            \textbf{\$0 USD} &
            \textbf{No requiere licencias ni infraestructura dedicada adicional} \\
        \midrule
    \end{tabular}
\end{table}

El costo incremental de licenciamiento e infraestructura de VulnSeeker es de \$0~USD, lo cual valida la
hipótesis de viabilidad económica planteada en el Capítulo~I y confirma que
la herramienta elimina las barreras financieras que actualmente impiden a
desarrolladores y PYMEs acceder a auditorías de seguridad automatizadas.

% ===========================================================================
% REFERENCIAS BIBLIOGRÁFICAS
% ===========================================================================
\newpage
\printbibliography[title=Referencias Bibliográficas]

\end{document}
